\chapter{Pulmonology}

\medterm{Atypical Pneumonia}-- Pneumonia with a gradual onset, relatively mild chest findings, and prominent systemic or extrapulmonary symptoms. Common causes include Mycoplasma pneumoniae, Chlamydophila pneumoniae, and respiratory viruses.

\medterm{Intercostal Retractions}-- Inward movement of the skin between the ribs during inspiration, indicating increased work of breathing. It is commonly seen in infants and children with airway obstruction, bronchiolitis, asthma, pneumonia, or respiratory failure.

\medterm{Pulmonary Arteriovenous Fistula}-- An abnormal direct connection between a pulmonary artery and vein that bypasses the capillary bed. It can cause hypoxemia, cyanosis, dyspnea, paradoxical emboli, or brain abscess.

\medterm{Silicosis}-- Fibrotic lung disease caused by inhalation of crystalline silica dust, typically in mining, quarrying, construction, or stone-cutting work. It causes progressive breathlessness and increases the risk of tuberculosis, chronic airflow limitation, and lung cancer.

\medterm{Acute Cough}-- Cough lasting < 3 weeks. Most common causes: URI, acute bronchitis, pneumonia. Usually self-limited.

\medterm{Acute Respiratory Distress Syndrome}-- ARDS: a syndrome of acute, diffuse inflammatory lung injury causing severe hypoxemic respiratory failure with bilateral infiltrates not due to cardiogenic pulmonary edema (PaO2/FiO2 ratio <=300). Causes: sepsis (most common), pneumonia, trauma, aspiration, pancreatitis, massive transfusion. Pathology: endothelial/alveolar injury, protein-rich edema, surfactant dysfunction, hyaline membranes; three phases—exudative, proliferative, fibrotic. Features: severe dyspnea, hypoxemia, refractory hypoxemia; diagnosis by Berlin criteria. Management: lung-protective ventilation (low tidal volume 6 mL/kg, plateau pressure <=30), PEEP, prone positioning, conservative fluids, neuromuscular blockade; steroids controversial. Mortality ~30-40 percent.

\medterm{Airway Clearance Techniques}-- Physical therapies to mobilize and remove secretions from the airways. Techniques include postural drainage, percussion, vibration, active cycle of breathing, autogenic drainage, positive expiratory pressure, oscillating PEP, and high-frequency chest wall oscillation. They are essential in CF, bronchiectasis, and neuromuscular disease.

\medterm{Allergen Immunotherapy}-- Gradual desensitization to specific allergens through repeated subcutaneous or sublingual exposure. It reduces allergic asthma symptoms and medication needs. Treatment requires 3--5 years for full benefit. It is indicated for allergic asthma with identified specific allergens not controlled by standard therapy.

\medterm{Allergic Bronchopulmonary Aspergillosis}-- A hypersensitivity reaction to Aspergillus fumigatus in asthma or CF patients. Causes central bronchiectasis, elevated total IgE, and eosinophilia. Treatment includes corticosteroids and itraconazole.

\medterm{Alpha-1 Antitrypsin Deficiency}-- An autosomal codominant genetic disorder caused by mutations in the SERPINA1 gene. It predisposes to early-onset panlobular emphysema (lower lobes), liver disease, and panniculitis. PiZZ genotype is most severe. Treatment includes AAT augmentation therapy and smoking avoidance.

\medterm{Anti-IgE Therapy}-- Omalizumab is a monoclonal antibody that binds free IgE, preventing it from binding to mast cells and basophils. It reduces allergic inflammation and is approved for moderate-to-severe persistent allergic asthma. It reduces exacerbations and the need for oral corticosteroids. It is administered by subcutaneous injection every 2--4 weeks.

\medterm{Anti-IL4 Receptor Therapy}-- Dupilumab is a monoclonal antibody blocking the IL-4 receptor alpha subunit, inhibiting both IL-4 and IL-13 signaling. It is approved for moderate-to-severe asthma with eosinophilic phenotype or oral corticosteroid dependence. It reduces exacerbations and improves lung function.

\medterm{Anti-IL5 Therapy}-- Monoclonal antibodies targeting interleukin-5, a key cytokine in eosinophil production and survival. Mepolizumab, reslizumab, and benralizumab (anti-IL-5 receptor) are approved for severe eosinophilic asthma. They reduce exacerbations and oral corticosteroid requirements in patients with blood eosinophils >= 300.

\medterm{ARDS}-- Acute respiratory distress syndrome; see Acute Respiratory Distress Syndrome. Acute diffuse inflammatory lung injury causing severe hypoxemia (PaO2/FiO2 <=300), bilateral infiltrates, and non-cardiogenic pulmonary edema. Common causes: sepsis, pneumonia, trauma, aspiration. Management: lung-protective ventilation (6 mL/kg tidal volume, low plateau pressure), PEEP, prone positioning, conservative fluid strategy, and treatment of the underlying cause. Complications: barotrauma, VAP, multi-organ failure; mortality 30-40 percent.

\medterm{Asbestosis}-- See public\_health.

\medterm{Aspirin-Exacerbated Respiratory Disease}-- A triad of asthma, nasal polyps, and aspirin sensitivity due to COX-1 inhibition. It involves overproduction of cysteinyl leukotrienes. Leukotriene receptor antagonists are specifically effective.

\medterm{Asthma}-- A chronic inflammatory disorder of the airways characterized by recurrent episodes of wheezing, breathlessness, chest tightness, and cough. Airflow limitation is often reversible. It involves eosinophilic inflammation, mucus hypersecretion, and smooth muscle hypertrophy.

\medterm{Asthma-COPD Overlap}-- Features of both asthma and COPD in the same patient. Characterized by persistent airflow limitation with features of both diseases: airway inflammation (eosinophilic and neutrophilic), bronchial hyperresponsiveness, and partially reversible obstruction. Treatment: inhaled corticosteroids plus long-acting bronchodilators. Distinguished from pure asthma (usually atopic, reversible) and pure COPD (smoking history, fixed obstruction).

\medterm{Brittle Asthma}-- Rare phenotype of severe asthma characterized by wide, chaotic, unpredictable peak expiratory flow variability ($> 40\%$ diurnal variation over weeks), classified as type 1 (chronic, persistent chaotic variability maintained over months despite maximal therapy) and type 2 (sudden onset of severe, life-threatening attacks against a background of apparently well-controlled asthma with normal pulmonary function between attacks; can be triggered by allergens, stress, or anaphylaxis). Most cases have type 1 brittle asthma; risk factors include atopy, food sensitivities, psychosocial factors, female sex, and severe airway hyperresponsiveness. Management: supervised high-dose inhaled corticosteroids and long-acting $\beta_2$-agonists; consider monoclonal antibodies (omalizumab, mepolizumab, benralizumab, dupilumab) for type 2-mediated disease; rapid-action bronchodilators with written action plan; autoinjector epinephrine for those with anaphylaxis comorbidity; consider subcutaneous terbutaline infusion or nocturnal noninvasive ventilation for refractory cases. High mortality risk if mismanaged.

\medterm{Bagassosis}-- An occupational hypersensitivity pneumonitis (extrinsic allergic alveolitis) caused by inhalation of dust from moldy bagasse—the fibrous residue of sugarcane after extraction. Caused by thermophilic actinomycetes (Saccharopolyspora rectivirgula) that colonize stored bagasse. Presents with fever, cough, dyspnea, and malaise 4--8 hours after exposure. Acute episodes resolve with corticosteroid treatment and avoidance; chronic exposure leads to pulmonary fibrosis.

\medterm{Berylliosis}-- A granulomatous lung disease from beryllium exposure. Chronic berylliosis resembles sarcoidosis. Diagnosis requires beryllium lymphocyte proliferation test.

\medterm{BiPAP}-- Bilevel Positive Airway Pressure, a form of non-invasive ventilation delivering two levels of pressure: IPAP (inspiratory) for ventilation and EPAP (expiratory) for alveolar recruitment. It is particularly effective for COPD exacerbations with hypercapnic respiratory failure. It reduces work of breathing and improves gas exchange.

\medterm{Bird Fancier's Lung}-- HP from avian proteins in bird droppings or feathers. Common among pigeon breeders. Chronic exposure leads to progressive fibrosis.

\medterm{Blue Bloater}-- A clinical phenotype of chronic bronchitis in COPD. Patients present with chronic cough, sputum production, cyanosis, right heart failure, and airway-centric inflammation. They tend to be younger, have more hypercapnia, and less emphysema on imaging.

\medterm{Bradypnea}-- Abnormally slow respiratory rate (adult <12 breaths/min). Causes: opioid/benzodiazepine/sedative overdose, central nervous system depression (stroke, trauma, increased ICP), hypothyroidism, hypothermia, and metabolic (e.g., hypokalemia). May lead to hypoventilation, hypercapnia, and respiratory failure. Management: airway protection, ventilation support, reversal of depressants (naloxone, flumazenil), and treatment of the cause. Contrast with apneustic/Cheyne-Stokes patterns in specific CNS lesions.

\medterm{Bronchial Carcinoid}-- Pulmonary neuroendocrine tumor arising from Kulchitsky cells of the bronchial mucosa; subclassified as typical carcinoid (TC, low-grade, $< 3$ mitoses/2 mm$^2$, no necrosis, excellent prognosis with $> 90\%$ 5-year survival) and atypical carcinoid (AC, intermediate-grade, 2-10 mitoses or necrosis present, higher recurrence/ metastasis). Typical are 3-4$\times$ more common; patients typically young adults (40s), more common in women, non-smoking association. Presentations: central tumors cause cough, hemoptysis, wheeze, post-obstructive pneumonia; peripheral tumors typically asymptomatic. Cushing syndrome from ectopic ACTH in $\sim 2\%$ of typical, up to $40-50\%$ of atypical. Carcinoid syndrome rare unless hepatic metastases. Diagnosis: CT chest, bronchoscopy with biopsy (cautious due to hemorrhage risk); octreotide scan/SSTR PET for staging. Treatment: anatomical surgical resection (lobe or segment) with lymphadenectomy; bronchial-sparing sleeve resection for central typical; radiation, chemotherapy, and somatostatin analogs for metastatic disease.

\medterm{Bronchial Biopsy}-- Tissue sampling of the bronchial wall during bronchoscopy. It is used to diagnose inflammatory conditions (sarcoidosis, eosinophilic bronchitis), infections, and malignancy. Endobronchial biopsy of visible lesions and transbronchial biopsy of parenchymal lesions provide complementary diagnostic information.

\medterm{Bronchiectasis}-- Permanent, abnormal dilation of bronchi caused by destruction of the bronchial wall from chronic airway inflammation and infection. Causes include CF, immune deficiency, primary ciliary dyskinesia, and post-infectious damage. Features include chronic cough, copious purulent sputum, and recurrent hemoptysis.

\medterm{Bronchitis}-- Inflammation of the bronchial mucosa. Acute bronchitis is typically viral, presents with cough (often productive), and resolves spontaneously. Cough with sputum production for at least 3 months in 2 consecutive years, often due to smoking. Treatment: supportive for acute; bronchodilators, smoking cessation, and pulmonary rehabilitation for chronic.

\medterm{Bronchoalveolar Lavage}-- Diagnostic procedure during bronchoscopy where saline is instilled and aspirated. Analyzed for cell counts, microbiology, and cytology. Useful for Pneumocystis pneumonia, pulmonary alveolar proteinosis, and sarcoidosis diagnosis.

\medterm{Bronchodilator Reversibility}-- Improvement in FEV1 > 12\% and > 200 mL after short-acting bronchodilator. Supports asthma diagnosis. Some COPD patients show partial reversibility.

\medterm{Bronchopleural Fistula}-- An abnormal communication between the bronchial tree and pleural space. It is a complication of pneumonectomy, lung surgery, or empyema. It causes persistent air leak, subcutaneous emphysema, and respiratory compromise. Treatment includes chest tube drainage, bronchoscopic intervention, and surgical repair.

\medterm{Centriacinar Emphysema}-- A subtype of emphysema in which destruction begins in the respiratory bronchioles and extends peripherally. It predominantly affects the upper lobes and is strongly associated with cigarette smoking. It is the most common form of emphysema and often coexists with chronic bronchitis.

\medterm{Chest Physiotherapy}-- Techniques to mobilize and clear secretions from the lungs, including postural drainage, percussion, vibration, and huffing. It is used in bronchiectasis, CF, and post-operative patients. Positive expiratory pressure devices and oscillating PEP devices are also used.

\medterm{Chronic Bronchitis}-- Defined clinically as a productive cough for at least three months in each of two consecutive years, after excluding other causes. It is characterized by hypertrophy of mucus-secreting glands in the bronchi (Reid index > 0.5), goblet cell hyperplasia, and chronic airway inflammation.

\medterm{Chronic Cough}-- Cough lasting > 8 weeks. Top 3 causes in non-smokers with normal CXR: upper airway cough syndrome, asthma, GERD. ACE inhibitor use is another common cause.

\medterm{Chronic Obstructive Pulmonary Disease}-- COPD is a preventable lung disease characterized by persistent respiratory symptoms and airflow limitation due to airway and/or alveolar abnormalities, usually caused by significant exposure to noxious particles or gases (primarily tobacco smoke). Includes emphysema and chronic bronchitis. Diagnosis: spirometry with post-bronchodilator FEV1/FVC <0.70. Management: smoking cessation, bronchodilators, inhaled corticosteroids, pulmonary rehabilitation, oxygen therapy, and lung volume reduction.

\medterm{Chronic Obstructive Pulmonary Disease (COPD)}-- COPD is a progressive lung disease characterized by persistent airflow limitation due to chronic bronchitis and emphysema, most commonly caused by cigarette smoking. Symptoms include dyspnea, chronic cough, sputum production, and frequent respiratory infections. Treatment involves bronchodilators, inhaled corticosteroids, pulmonary rehabilitation, supplemental oxygen, and smoking cessation to reduce exacerbations and improve quality of life.

\medterm{Coal Worker's Pneumoconiosis}-- Pneumoconiosis from coal dust. Simple CWP shows upper lobe macules. Complicated CWP (progressive massive fibrosis) involves coalescent macules. May be associated with Caplan syndrome.

\medterm{Coarse Crackles}-- Lower-pitched, longer-duration crackles heard throughout inspiration, caused by fluid or secretions in larger airways. They are heard in COPD, bronchiectasis, pneumonia, and pulmonary edema. They are less specific than fine crackles.

\medterm{Combination Inhalers}-- Inhalers containing both an inhaled corticosteroid and a long-acting beta-2 agonist. Examples include fluticasone/salmeterol (Advair), budesonide/formoterol (Symbicort), and fluticasone furoate/vilanterol (Breo Ellipta). They simplify asthma and COPD maintenance therapy and improve adherence.

\medterm{Crackles (Rales)}-- Discontinuous, non-musical adventitious lung sounds produced by the sudden explosive opening of collapsed small airways and alveoli during inspiration or the bubbling of air through fluid in the bronchioles. Fine crackles: high-pitched, short duration, late inspiratory; characteristic of interstitial lung disease (idiopathic pulmonary fibrosis) and early congestive heart failure. Coarse crackles: low-pitched, moist, early inspiratory/expiratory; characteristic of pulmonary edema, pneumonia, bronchitis, and bronchiectasis.

\medterm{Cryobiopsy}-- A bronchoscopic technique using rapid freezing to obtain larger lung tissue samples than conventional transbronchial biopsy. It provides better histopathological specimens for diagnosing interstitial lung diseases. It requires general anesthesia and bronchoscopic expertise. It is less invasive than surgical lung biopsy.

\medterm{Dead Space Ventilation}-- Ventilation of alveoli without perfusion (V/Q = infinity). Anatomic dead space is the conducting airways (~150 mL). Physiologic dead space includes alveoli without perfusion. Increased in PE, hypotension, and emphysema. Measured by the Bohr equation.

\medterm{Desquamative Interstitial Pneumonia}-- A rare ILD characterized by uniform filling of alveoli with macrophages containing brown pigment (hemosiderin-laden). It is strongly associated with smoking. HRCT shows bilateral ground-glass opacities with lower lobe predominance. Smoking cessation and corticosteroids are the main treatments.

\medterm{Diffusing Capacity (DLCO)}-- Ability of lungs to transfer gas from alveoli to blood. Reduced in emphysema, ILD, and anemia. Increased in pulmonary hemorrhage and left-to-right shunts.

\medterm{Digital Clubbing}-- Bulbous enlargement of distal phalanges with loss of nail-bed angle (> 180 degrees). Associated with lung cancer, ILD, bronchiectasis, and cyanotic heart disease.

\medterm{Dyspnoea}-- The subjective sensation of difficult or labored breathing, also known as shortness of breath (SOB). It is a symptom common to many cardiopulmonary and systemic conditions. Causes: respiratory (asthma, COPD, pneumonia, pulmonary embolism, pneumothorax, interstitial lung disease), cardiac (heart failure, myocardial ischemia, arrhythmia, valvular disease), and others (anemia, deconditioning, anxiety, metabolic acidosis, thyrotoxicosis). Assessed using the Borg Scale or modified Medical Research Council (mMRC) scale. Acute dyspnea is a medical emergency requiring rapid assessment of airway, breathing, and circulation; immediate oxygen supplementation and treatment of the underlying cause (bronchodilators, diuretics, anticoagulation, etc.). Chronic dyspnea requires systematic evaluation with pulmonary function testing, echocardiography, and imaging.

\medterm{Emphysema}-- A form of COPD characterized by permanent, abnormal enlargement of airspaces distal to the terminal bronchioles, accompanied by destruction of their walls. Centriacinar emphysema affects the respiratory bronchioles and is associated with smoking. Panlobular emphysema involves the entire acinus and is associated with alpha-1 antitrypsin deficiency.

\medterm{Endobronchial Foreign Body}-- A foreign object lodged in a bronchus, causing partial or complete airway obstruction. It may cause localized wheezing, atelectasis, and post-obstructive pneumonia. It is most common in children and may be initially misdiagnosed as asthma. Bronchoscopy is both diagnostic and therapeutic.

\medterm{Endobronchial Valve}-- A one-way valve placed bronchoscopically in a target airway to achieve lobar atelectasis in emphysema. It allows air and secretions to exit but prevents air from entering. It is used for bronchoscopic lung volume reduction in selected patients with heterogeneous emphysema and intact fissures.

\medterm{Exacerbation of COPD}-- An acute worsening of respiratory symptoms beyond normal day-to-day variations, requiring change in medication. It is commonly triggered by respiratory infections (viral, bacterial), air pollution, or non-compliance. Treatment includes increased bronchodilators, systemic corticosteroids, antibiotics if indicated, and oxygen therapy.

\medterm{Exercise-Induced Bronchoconstriction}-- Airway narrowing during or after exercise, presenting with cough, wheeze, and dyspnea. It occurs in up to 90\% of asthmatics and is triggered by heat and water loss from the airways. Prevention includes pre-exercise SABA.

\medterm{Exhaled Nitric Oxide}-- A non-invasive measure of airway eosinophilic inflammation. FeNO > 25 ppb in adults suggests eosinophilic airway inflammation and predicts response to inhaled corticosteroids. It is useful for diagnosing eosinophilic asthma, monitoring treatment response, and differentiating asthma from other causes of cough.

\medterm{Expectoration}-- The act of coughing up and spitting out mucus, sputum, or phlegm from the respiratory tract; sputum production. Sputum characteristics aid diagnosis: purulent (bacterial infection), frothy pink (pulmonary edema), rust-colored (pneumococcal pneumonia), green (Pseudomonas, COPD exacerbation), bloody (hemoptysis—TB, malignancy, bronchiectasis), mucoid (asthma). Assessment includes quantity, color, viscosity, and microbiologic studies (Gram stain, culture). Treatment targets the underlying cause; mucolytics/expectorants and hydration may help.

\medterm{FEV1/FVC Ratio}-- Fraction of air exhaled in first second relative to total. < 0.70 post-bronchodilator indicates obstruction. Normal or elevated in restrictive disease.

\medterm{Fine Crackles}-- High-pitched, brief, discontinuous sounds heard at the end of inspiration, caused by opening of small airways. They are characteristic of ILD, pulmonary fibrosis, and heart failure. They are also called Velcro crackles due to their resemblance to the sound of Velcro being pulled apart.

\medterm{Flutter Valve}-- A handheld device combining PEP therapy with oscillations during exhalation. The oscillations create shear forces that loosen mucus from airway walls, while the PEP keeps airways open. It is used for airway clearance in bronchiectasis, CF, and COPD.

\medterm{Forced Oscillation Technique}-- Similar to impulse oscillometry, applying small pressure oscillations at various frequencies during tidal breathing to measure respiratory resistance and reactance. It provides information about airway mechanics without requiring forced expiratory maneuvers. It is useful for detecting small airway dysfunction and monitoring treatment response in asthma and COPD.

\medterm{Forced Vital Capacity}-- Total air exhaled after maximal inspiration. Reduced in both obstructive and restrictive disease. In restrictive disease, FVC reduced with preserved ratio.

\medterm{Haemoptysis}-- Coughing up blood or blood-stained sputum from the lower respiratory tract. See \hyperlink{hematology}{hematology}.

\medterm{Hemoptysis}-- Expectoration of blood from the lower respiratory tract. Must be differentiated from hematemesis and epistaxis. Causes: bronchitis, pneumonia, TB, bronchiectasis, lung cancer, PE.

\medterm{Hypersensitivity Pneumonitis}-- An immune-mediated lung disease from environmental antigen inhalation (molds, bird proteins). Acute, subacute, and chronic forms exist. HRCT shows ground-glass opacities, air trapping, and centrilobular nodules. Chronic HP progresses to fibrosis.

\medterm{Hypertonic Saline Inhalation}-- Nebulized 3--7\% sodium chloride solution that draws water into the airway surface liquid by osmosis, improving mucociliary clearance. It is used in CF and bronchiectasis. It may cause bronchospasm and should be preceded by a bronchodilator. It reduces exacerbations in CF.

\medterm{Hypoxia}-- Reduced oxygen availability at the tissue level, defined as arterial partial pressure of oxygen (PaO$_2$) <60 mmHg or arterial oxygen saturation (SpO$_2$) <90\% on room air. Five types: (1) Hypoxic hypoxia (low arterial PO$_2$)—caused by low inspired oxygen (high altitude), hypoventilation, V/Q mismatch, shunt, or diffusion impairment (pneumonia, COPD, pulmonary edema, ARDS). (2) Anaemic hypoxia (reduced oxygen-carrying capacity)—anemia, carbon monoxide poisoning, methaemoglobinaemia. (3) Stagnant/hypoperfusion hypoxia (reduced blood flow)—heart failure, shock, peripheral vascular disease. (4) Histotoxic hypoxia (impaired tissue utilisation)—cyanide poisoning (blocks cytochrome c oxidase), sepsis. (5) Demand hypoxia (increased O$_2$ requirement)—exercise, hypermetabolic states (thyrotoxicosis, malignant hyperthermia). Assessment: ABG, pulse oximetry, and calculation of alveolar-arterial (A-a) gradient. Consequences: cellular anaerobic metabolism, lactic acidosis, ATP depletion, and cell death. Treatment: supplemental oxygen (high-flow nasal cannula, non-invasive ventilation, mechanical ventilation), treat underlying cause. The hypoxic ventilatory response is mediated by peripheral chemoreceptors (carotid and aortic bodies).

\medterm{Idiopathic Pulmonary Fibrosis}-- A chronic, progressive fibrosing interstitial pneumonia of unknown cause. Histopathology shows usual interstitial pneumonia (UIP). Features include exertional dyspnea, dry cough, bibasilar crackles, and digital clubbing. Median survival is 3--5 years.

\medterm{Impulse Oscillometry}-- A pulmonary function testing technique that measures respiratory impedance by applying small pressure oscillations during tidal breathing. It assesses large airway resistance (R5), small airway resistance (R5-R20), and reactance. It is more sensitive than spirometry for detecting early small airway disease and is useful in patients who cannot perform spirometry (young children, elderly).

\medterm{Incentive Spirometry}-- A device that encourages deep breathing to prevent atelectasis after surgery. The patient inhales slowly and deeply through the device, maintaining maximal inspiration for several seconds. It promotes alveolar recruitment and prevents mucus plugging. It is standard post-operative care.

\medterm{Induced Sputum Analysis}-- A non-invasive method of assessing airway inflammation by examining sputum cell counts and differentials obtained after hypertonic saline inhalation. Sputum eosinophils > 3\% indicate eosinophilic airway inflammation (asthma, EGPA). Neutrophilic inflammation suggests infection, COPD, or bronchiectasis. It guides anti-inflammatory therapy decisions.

\medterm{Inhaled Corticosteroids}-- Anti-inflammatory medications delivered directly to the airways, reducing airway inflammation, mucus production, and bronchial hyperresponsiveness. Examples include fluticasone, budesonide, beclomethasone, and mometasone. They are the cornerstone of asthma maintenance therapy. Side effects include oral candidiasis and dysphonia.

\medterm{Inspiratory Crackles}-- Fine, discontinuous sounds heard during inspiration, caused by the sudden opening of collapsed airways. They are a hallmark of interstitial lung disease (bibasilar, Velcro-like crackles) and pulmonary fibrosis. They are also heard in pneumonia, atelectasis, and heart failure.

\medterm{Interstitial Lung Disease}-- A heterogeneous group of disorders with inflammation and fibrosis of lung parenchyma, causing restrictive physiology (reduced FVC and TLC with normal or elevated FEV1/FVC ratio).

\medterm{Interstitial Pneumonia}-- Inflammation centered on the interstitium. Caused by Mycoplasma, Chlamydophila, Legionella, and viruses. Presents with dry cough and diffuse crackles.

\medterm{Ipratropium Bromide}-- An inhaled anticholinergic (antimuscarinic) bronchodilator that blocks acetylcholine at airway smooth muscle muscarinic receptors. It reduces bronchospasm and mucus secretion. It is used in COPD and acute asthma exacerbations. It is less effective than beta-2 agonists in asthma but more effective in COPD.

\medterm{Kartagener Syndrome}-- A subset of primary ciliary dyskinesia with situs inversus, chronic sinusitis, and bronchiectasis. Caused by dynein defects impairing both mucociliary clearance and embryonic nodal cilia rotation.

\medterm{Kartagener Triad}-- The classic triad of situs inversus, chronic sinusitis, and bronchiectasis, occurring in approximately 50\% of patients with primary ciliary dyskinesia. It is caused by dynein defects impairing both mucociliary clearance and embryonic nodal cilia rotation.

\medterm{Lfgren Syndrome}-- Acute sarcoidosis with bilateral hilar lymphadenopathy, erythema nodosum, and ankle polyarthralgia. Excellent prognosis with spontaneous resolution. Associated with HLA-DR3.

\medterm{Long-Acting Beta-2 Agonists}-- Bronchodilators that provide sustained airway smooth muscle relaxation for 12--24 hours. Examples include salmeterol, formoterol, indacaterol, and olodaterol. They are used for maintenance therapy in asthma and COPD, always in combination with inhaled corticosteroids in asthma.

\medterm{Lower Airway Obstruction}-- Narrowing of airways below the carina, causing expiratory wheezing and airflow obstruction on spirometry. Causes include asthma, COPD, bronchiectasis, and foreign body. It is characterized by scooped-out expiratory limb on flow-volume loops and reduced FEV1/FVC ratio.

\medterm{Lung Disease}-- Lung disease encompasses disorders affecting the airways, parenchyma, pleura, pulmonary vasculature, and chest wall. Major categories: obstructive (asthma, COPD—chronic bronchitis and emphysema), restrictive (pulmonary fibrosis, sarcoidosis, asbestosis), infectious (pneumonia, tuberculosis), pulmonary vascular (pulmonary embolism, pulmonary arterial hypertension), and neoplastic (lung cancer, the leading cause of cancer death). Common symptoms: dyspnea, cough (dry or productive), wheezing, hemoptysis, chest pain, and weight loss. Diagnosis: pulmonary function tests (spirometry, DLCO), chest imaging (X-ray, CT), bronchoscopy, and biopsy. Management: smoking cessation (most effective intervention), bronchodilators (beta-agonists, antimuscarinics), inhaled corticosteroids, oxygen therapy, pulmonary rehabilitation, antimicrobials for infection, and lung transplantation for end-stage disease.

\medterm{Lung Transplant}-- Surgical replacement of one or both lungs from a donor for end-stage lung disease: COPD (most common), idiopathic pulmonary fibrosis, cystic fibrosis, pulmonary hypertension, and other advanced conditions. Types: single, bilateral, heart-lung. Evaluation: multidisciplinary, age, comorbidities, functional status, infection; waiting list allocation by severity. Immunosuppression (calcineurin inhibitors, mycophenolate, steroids) prevents rejection. Complications: primary graft dysfunction, acute/chronic rejection (bronchiolitis obliterans), infection (CMV, fungal), airway anastomotic complications, and malignancy. Survival ~50-90 percent at 5 years by indication.

\medterm{Lung Transplantation}-- Surgical replacement of a diseased lung with a healthy lung from a donor. Indications include IPF, COPD, CF, PAH, and bronchiectasis. Options include single lung transplant, double lung transplant, or heart-lung transplant. Post-transplant immunosuppression is required lifelong. Complications include rejection and infection.

\medterm{Lymphocytic Interstitial Pneumonia}-- A rare ILD characterized by diffuse polyclonal lymphocytic infiltration of the lung parenchyma. It is strongly associated with Sjgren syndrome and HIV. HRCT shows ground-glass opacities, thin-walled cysts, and lymphadenopathy. Treatment is corticosteroids for symptomatic disease.

\medterm{Metered-Dose Inhaler}-- A pressurized inhaler delivering a metered dose of medication. Requires proper technique for effective delivery. Spacer devices improve drug delivery and reduce oropharyngeal deposition. Most common delivery device for asthma and COPD maintenance.

\medterm{Methacholine Challenge Test}-- See medical\_tests.

\medterm{Middle Lobe Syndrome}-- Chronic or recurrent atelectasis and infection of the right middle lobe due to bronchial obstruction. The long, narrow right middle lobe bronchus is susceptible to compression by lymph nodes, tumors, or mucous plugging. It may be caused by bronchiectasis, malignancy, or granulomatous disease. Treatment depends on the underlying cause.

\medterm{Mucolytic Therapy}-- Medications that reduce mucus viscosity and elasticity, facilitating clearance. N-acetylcysteine (NAC) breaks disulfide bonds in mucus glycoproteins. Dornase alfa (Pulmozyme) cleaves extracellular DNA in CF sputum. Hypertonic saline aerosol draws water into the airway surface liquid. They are used in CF, bronchiectasis, and chronic bronchitis.

\medterm{Mucus Plugging}-- Accumulation of thick, tenacious mucus within airways causing obstruction. It is common in asthma, CF, bronchiectasis, and chronic bronchitis. It impairs gas exchange, promotes infection, and causes atelectasis. Treatment includes airway clearance techniques, mucolytics, and bronchodilators.

\medterm{Nebulizer}-- A device that converts liquid medications into a mist for inhalation. Used for bronchodilators (albuterol, ipratropium) and corticosteroids in acute exacerbations. More effective than MDI in patients with severe obstruction.

\medterm{Non-Asthmatic Eosinophilic Bronchitis}-- Chronic cough with airway eosinophilia (induced sputum eosinophils > 3\%) but without bronchial hyperresponsiveness. It is a common cause of chronic cough and does not respond to bronchodilators. It responds well to inhaled corticosteroids. It is distinct from asthma in the absence of variable airflow obstruction.

\medterm{Non-Invasive Ventilation}-- Positive pressure ventilation delivered via a face mask or nasal mask without endotracheal intubation. BiPAP (bilevel positive airway pressure) is most commonly used for acute hypercapnic respiratory failure (COPD exacerbation) and cardiogenic pulmonary edema. It reduces the need for intubation and improves outcomes in selected patients.

\medterm{Nonspecific Interstitial Pneumonia}-- A pattern with uniform, temporally homogeneous inflammation and/or fibrosis. Most common ILD associated with connective tissue disease. HRCT shows bilateral ground-glass opacities with subpleural sparing.

\medterm{Occupational Asthma}-- Asthma from workplace exposures to sensitizers (isocyanates, flour, latex) or irritants. Improves away from work and worsens at work. Serial peak flow monitoring supports diagnosis.

\medterm{Occupational COPD}-- COPD from workplace dust, chemical, or fume exposure independent of smoking. Accounts for 15--20\% of COPD cases. High-risk occupations: mining, construction, farming.

\medterm{Orthopnoea}-- Shortness of breath that occurs when lying flat and is relieved by sitting or standing upright. Orthopnoea is a classic symptom of left-sided heart failure (elevated pulmonary venous pressure worsens pulmonary congestion in the supine position). Other causes: severe COPD, asthma exacerbation, diaphragmatic paralysis, ascites, and obesity. It is quantified by the number of pillows required to sleep comfortably (e.g., two-pillow orthopnoea). Pathophysiology: supine position increases venous return (preload) and reduces lung volumes (cephalad displacement of the diaphragm and abdominal contents), worsening V/Q mismatch and increasing pulmonary capillary pressure. Orthopnoea may develop acutely in pulmonary edema (cardiac asthma) or gradually in worsening chronic heart failure. Related: paroxysmal nocturnal dyspnoea (PND—sudden breathlessness during sleep, relieved by sitting up, also CHF). Management: treat the underlying cause—diuretics, vasodilators, and optimisation of heart failure therapy.

\medterm{Oxygen Saturation}-- The percentage of hemoglobin saturated with oxygen in arterial blood; normal 95-100 percent. Measured noninvasively by pulse oximetry (SpO2) and invasively by co-oximetry (SaO2). Target levels vary by context: >94 percent general; 88-92 percent in patients with hypercapnic respiratory failure (COPD); permissive targets in ARDS. Reductions indicate hypoxemia from ventilation-perfusion mismatch, shunt, hypoventilation, or reduced inspired oxygen. False readings from poor perfusion, anemia, carboxyhemoglobin. Guides oxygen therapy.

\medterm{Panlobular Emphysema}-- A subtype involving diffuse destruction of the entire acinus. It predominantly affects the lower lobes and is classically associated with alpha-1 antitrypsin deficiency. The destruction is more uniform compared to centriacinar emphysema.

\medterm{Paraseptal Emphysema}-- A form of emphysema characterized by destruction of alveoli adjacent to the pleura. It typically affects young adults and predisposes to spontaneous pneumothorax due to rupture of subpleural blebs. It may be associated with Marfan syndrome.

\medterm{Peak Expiratory Flow Rate (PEFR)}-- The maximum flow rate achieved during a forced expiratory manoeuvre after maximal inspiration, measured in litres per minute (L/min) using a peak flow meter. PEFR reflects the degree of airflow obstruction and is effort-dependent. Normal values are predicted based on age, sex, and height (e.g., healthy adult male 600 L/min, female 400 L/min). Uses: diagnosis and monitoring of asthma (diurnal variability >20\% supports asthma diagnosis), assessment of exacerbation severity (PEFR <50\% of predicted indicates severe exacerbation—NICE/BTS guidelines), and guiding asthma action plans (personal best, zone system). PEFR <33\% of predicted is life-threatening. Limitations: effort-dependent, poor technique, less reproducible than FEV$_1$.

\medterm{Percussion Note}-- The sound produced by percussing the chest wall. Normal lung produces a resonant note. Dullness suggests consolidation or fluid. Hyperresonance suggests pneumothorax or emphysema. Stony dullness suggests pleural effusion. Tympany suggests large pneumothorax.

\medterm{Persistent Asthma}-- Asthma requiring daily controller therapy. Mild persistent: symptoms > 2 days/week but not daily. Moderate persistent: daily symptoms, nighttime awakenings > 1/week. Severe persistent: symptoms throughout the day, frequent nighttime awakenings, SABA use several times/day, extremely limited activity.

\medterm{Pink Puffer}-- A clinical phenotype of emphysema in COPD. Patients present with severe dyspnea, minimal cough, thin habitus, pursed-lip breathing, accessory muscle use, and relatively normal PaO2 with late hypercapnia. They tend to be older and have more extensive emphysema.

\medterm{Pneumoconiosis}-- ILDs from mineral dust inhalation (silicosis, asbestosis, coal worker's, berylliosis). They share interstitial fibrosis and risk of infection and malignancy.

\medterm{Pneumomediastinum}-- Air in the mediastinum, often from alveolar rupture with air tracking along bronchovascular sheaths. Causes include asthma exacerbation, trauma, Valsalva maneuver, and esophageal rupture (Boerhaave syndrome). It presents with chest pain, subcutaneous emphysema, and Hamman's crunch (mediastinal crunch synchronous with heartbeat). Usually self-limited.

\medterm{Positive Pressure Ventilation}-- Mechanical ventilation that delivers air/oxygen under positive pressure to support or replace spontaneous breathing. Noninvasive (NIV): CPAP (continuous) and BiPAP (two levels) via mask for obstructive sleep apnea, cardiogenic pulmonary edema, and hypercapnic COPD. Invasive: endotracheal tube/tracheostomy delivering volume- or pressure-controlled breaths. Modes: assist-control, SIMV, pressure support, PRVC. Benefits: improved oxygenation, reduced work of breathing, alveolar recruitment. Complications: barotrauma/pneumothorax, VAP, hemodynamic compromise, hyperinflation. Requires monitoring and lung-protective settings.

\medterm{Post-Operative Atelectasis}-- Collapse of alveoli after surgery, the most common post-operative pulmonary complication. Risk factors include general anesthesia, abdominal surgery, obesity, smoking, and pre-existing lung disease. Prevention includes incentive spirometry, early ambulation, and adequate pain control.

\medterm{Primary Ciliary Dyskinesia}-- An autosomal recessive disorder of ciliary structure, most commonly caused by dynein arm defects. Leads to impaired mucociliary clearance, chronic sinopulmonary infections, bronchiectasis, and situs inversus (50\% of cases).

\medterm{Pulmonary Fibrosis}-- Scarring replacing normal alveolar architecture. Causes: IPF, connective tissue disease, asbestos, radiation, drugs (bleomycin, methotrexate). Features: progressive dyspnea, dry cough, digital clubbing.

\medterm{Pulse Oximetry}-- A noninvasive method measuring oxygen saturation by spectrophotometry: light at two wavelengths (red and infrared) passes through a pulsatile vascular bed (finger, ear, toe); absorption difference estimates oxyhemoglobin fraction (SpO2). Normal 95-100 percent. Used widely in monitoring (perioperative, emergency, sleep, home oxygen). Limitations: inaccurate with poor perfusion, motion, nail polish, anemia (despite adequate O2 content), and elevated carboxyhemoglobin (falsely normal in CO poisoning); does not assess ventilation (PaCO2).

\medterm{Respiratory Acidosis}-- Acid-base disturbance due to accumulation of CO2 from hypoventilation, lowering pH (acidemia). Causes: respiratory depression (opioids, sedatives, CNS disease), airway obstruction, neuromuscular weakness (GBS, myasthenia), chest wall restriction, severe COPD/asthma, and ventilation-perfusion mismatch. Compensation: renal bicarbonate retention (slow, days). Diagnosis: ABG—elevated PaCO2 (>45 mmHg), low pH, increased HCO3 if chronic. Management: treat cause, airway support, noninvasive/invasive ventilation; avoid over-correction; chronic compensation protects pH.

\medterm{Respiratory Alkalosis}-- Acid-base disturbance from excessive elimination of CO2 by hyperventilation, raising pH (alkalemia). Causes: anxiety/hyperventilation, hypoxia, pulmonary embolism, sepsis, fever, pain, salicylate toxicity, liver failure, pregnancy, high altitude. Compensation: renal bicarbonate excretion (days). Features: paresthesias, lightheadedness, tetany, carpopedal spasm (hypocalcemia symptoms). Diagnosis: ABG—low PaCO2 (<35 mmHg), high pH. Management: treat underlying cause, rebreathing/calming for anxiety; avoid over-sedation; check for causes like PE/ARDS which are dangerous.

\medterm{Respiratory Failure}-- Inadequate gas exchange resulting in hypoxemia (PaO2 <60 mmHg), hypercapnia (PaCO2 >50 mmHg), or both. Type 1 (hypoxemic): V/Q mismatch, shunt, diffusion impairment. Type 2 (hypercapnic): hypoventilation from neuromuscular disease, COPD exacerbation, or opioid overdose. Management: oxygen therapy, non-invasive ventilation, mechanical ventilation, and treatment of the underlying cause.

\medterm{Respiratory Health}-- The proper functioning of the lungs and airways, enabling efficient gas exchange and oxygen delivery to the body. Maintaining respiratory health involves avoiding tobacco smoke, minimizing exposure to air pollutants and occupational hazards, and receiving vaccinations against respiratory infections like influenza and pneumococcal disease. Regular exercise strengthens respiratory muscles and improves lung capacity, while prompt treatment of infections and management of chronic conditions such as asthma or COPD help preserve lung function over a lifetime.

\medterm{Respiratory Infections}-- Illnesses caused by pathogens including viruses, bacteria, and fungi that affect the airways and lungs. Common upper respiratory infections include the common cold, sinusitis, and pharyngitis, while lower respiratory infections such as bronchitis and pneumonia can be more serious, particularly in older adults and immunocompromised individuals. Treatment varies by cause, with viral infections often managed symptomatically and bacterial infections requiring antibiotics; prevention through hand hygiene and vaccination is key.

\medterm{Respiratory System}-- The anatomical and physiological system responsible for gas exchange—delivering oxygen to the blood and removing carbon dioxide. Components: (1) Upper respiratory tract—nose, nasal cavity, paranasal sinuses, pharynx (nasopharynx, oropharynx, laryngopharynx). The nose warms, humidifies, and filters inspired air; the nasal mucosa contains olfactory receptors. (2) Lower respiratory tract—larynx (voice production, airway protection—epiglottis, vocal cords), trachea (C-shaped cartilage rings, 10--12 cm long), bronchi (right main bronchus—wider, shorter, more vertical—common site of aspiration), bronchioles (no cartilage, smooth muscle—asthma), alveolar ducts, and alveoli (300--500 million, total surface area ~70--100 m\(^2\)), where gas exchange occurs across the thin alveolar-capillary membrane (0.2--0.5 \(\mu\)m). (3) Lungs—right lung (three lobes: upper, middle, lower) and left lung (two lobes: upper, lower—the cardiac notch accommodates the heart). The pleura (visceral and parietal) lines the lungs and chest wall; the pleural cavity contains a thin layer of surfactant-rich fluid for lubrication. (4) Chest wall and muscles of respiration—diaphragm (primary inspiratory muscle, innervated by phrenic nerve C3--C5), intercostal muscles (external—inspiration; internal—expiration), and accessory muscles (sternocleidomastoid, scalenes—used in respiratory distress). Ventilation: the rhythmic contraction of the diaphragm and intercostal muscles creates negative intrapleural pressure, drawing air into the lungs. Control: the respiratory center in the medulla oblongata generates the breathing rhythm, modulated by chemoreceptors (central—CO$_2$/pH; peripheral—O$_2$, CO$_2$, pH in carotid and aortic bodies).

\medterm{Respiratory Tract Infections}-- Infections affecting any part of the respiratory tract, classified by anatomical location into upper respiratory tract infections (URTIs) and lower respiratory tract infections (LRTIs). URTIs: common cold (rhinovirus, coronavirus, adenovirus), pharyngitis (Group A Streptococcus—bacterial, adenovirus, EBV), tonsillitis (Streptococcus pyogenes), sinusitis (Streptococcus pneumoniae, Haemophilus influenzae, Moraxella catarrhalis, viruses), otitis media (S. pneumoniae, H. influenzae), laryngitis (viruses, vocal strain). LRTIs: acute bronchitis (viruses—influenza, RSV; Bordetella pertussis), bronchiolitis (RSV—infants), pneumonia (community-acquired—S. pneumoniae, H. influenzae, Mycoplasma pneumoniae, Legionella, Chlamydia psittaci, viruses; hospital-acquired—Pseudomonas aeruginosa, MRSA, Enterobacteriaceae; aspiration—anaerobes). Risk factors: extremes of age, smoking, COPD, asthma, immunocompromise, cystic fibrosis, structural lung disease, and institutionalisation. Antibiotic resistance is a growing concern (MRSA, ESBL-producing Enterobacteriaceae, multidrug-resistant S. pneumoniae). Diagnosis: clinical (cough, sputum, fever, dyspnoea, chest pain), chest X-ray (consolidation in pneumonia), sputum culture, blood culture, nasopharyngeal swab (viral PCR), and urinary antigen testing (S. pneumoniae, Legionella). Prevention: hand hygiene, vaccination (influenza, pneumococcal, pertussis, RSV—maternal, palivizumab for high-risk infants), and smoking cessation.

\medterm{Rhonchi}-- Low-pitched, continuous adventitious sounds heard primarily during expiration, caused by secretions in the large airways. They may clear with coughing. They are heard in COPD, bronchitis, and bronchiectasis. They are different from wheezes in their lower pitch and coarser quality.

\medterm{Small Airway Disease}-- Disease affecting airways with diameter < 2 mm (bronchioles). It is the earliest site of damage in COPD and is often subclinical until advanced. It is characterized by narrowing, mucus plugging, and inflammation of small airways. PFTs may show air trapping (increased RV/TLC) before FEV1 decline.

\medterm{Spirometry}-- Most common PFT measuring air volume and flow during forced breathing. Key measurements: FVC, FEV1, FEV1/FVC ratio. Used to diagnose and classify obstructive and restrictive diseases.

\medterm{Status Asthmaticus}-- A severe, life-threatening asthma exacerbation unresponsive to standard bronchodilator therapy. Characterized by severe airflow obstruction, hypoxemia, and hypercapnia. Treatment includes continuous nebulized albuterol, IV magnesium sulfate, systemic corticosteroids, and possibly intubation.

\medterm{Subcutaneous Emphysema}-- Air in the subcutaneous tissues, often from pneumothorax, pneumomediastinum, or chest trauma. It causes a characteristic crepitus on palpation. It is usually benign and resolves when the underlying air leak is treated. Severe cases may cause skin necrosis and compartment syndrome.

\medterm{Supplemental Oxygen Therapy}-- Administration of oxygen above atmospheric concentration. Indicated for hypoxemia (PaO2 < 60 mmHg or SpO2 < 90\%). Delivered via nasal cannula (1--6 L/min), simple mask, non-rebreather mask, or mechanical ventilation. Target SpO2 88--92\% in COPD patients to avoid CO2 retention.

\medterm{Tachypnea}-- Abnormally rapid respiratory rate (adult >20 breaths/min). Physiologic (exercise, anxiety) or pathologic: hypoxia, fever, sepsis, metabolic acidosis (Kussmaul), pulmonary embolism, pneumonia, heart failure, pain, and compensation for metabolic derangement. Combined with other signs aids diagnosis; persistent tachypnea with distress suggests significant illness. Neonates and infants normally breathe faster. Evaluation: oxygen saturation, ABG, imaging, and treatment of the underlying cause.

\medterm{Tachypnoea}-- Abnormally rapid breathing, defined as a respiratory rate >20 breaths per minute in adults (age-dependent: >60 in neonates, >40 in 1--12 months, >30 in 1--5 years, >20 in >5 years). Tachypnoea is a non-specific sign of respiratory or metabolic distress. Causes: (1) Respiratory—hypoxia (pneumonia, asthma/COPD exacerbation, pulmonary embolism, pneumothorax, pulmonary edema, ARDS, interstitial lung disease), airway obstruction, and hypercapnia (COPD, respiratory muscle weakness). (2) Cardiac—heart failure (pulmonary congestion stimulates vagal irritant receptors). (3) Metabolic—metabolic acidosis (diabetic ketoacidosis—Kussmaul breathing, uraemia, lactic acidosis, salicylate poisoning), fever (respiratory rate increases 2--4 breaths/min per 1°C rise). (4) Other—anxiety (panic attack—psychogenic hyperventilation), pain, anemia, hyperthyroidism, pregnancy (progesterone stimulates respiratory center), exercise, and high altitude. (5) Drugs: salicylates, theophylline, $\beta$-agonists, and other stimulants. Tachypnoea is often accompanied by shallow breathing (low tidal volume). Assessment: respiratory rate, SpO$_2$, ABG, chest X-ray, ECG, and focused history. Persistent tachypnoea with hypoxia is a sign of serious illness requiring urgent investigation and treatment.

\medterm{Talcosis}-- Pneumoconiosis from talc dust inhalation or IV injection of crushed oral medications. Causes upper lobe nodular fibrosis. HRCT may show ``galaxy sign'' of calcified nodules.

\medterm{Thermoplasty}-- A bronchoscopic procedure that uses radiofrequency energy to ablate smooth muscle in the airways, reducing bronchospasm. It is approved for severe, persistent asthma in adults not controlled with ICS/LABA. It reduces exacerbations and emergency department visits but does not improve lung function.

\medterm{Tiotropium Bromide}-- A long-acting muscarinic antagonist (LAMA) bronchodilator used for maintenance therapy in COPD. Once-daily dosing provides sustained bronchodilation. It reduces exacerbations and improves lung function. It is often combined with a LABA for additive bronchodilation.

\medterm{Tracheal Stenosis}-- Narrowing of the trachea, often a complication of prolonged endotracheal intubation or tracheostomy. It presents with progressive dyspnea, stridor, and wheezing that may be mistaken for asthma. Diagnosis is by bronchoscopy or CT. Treatment is balloon dilation, stenting, or surgical resection.

\medterm{Transbronchial Biopsy}-- Bronchoscopic biopsy of lung parenchyma obtained through the bronchial wall. It is used to diagnose interstitial lung diseases (sarcoidosis, HP), reject lung transplant rejection, and sample peripheral lung lesions. It carries a higher risk of pneumothorax than endobronchial biopsy.

\medterm{Usual Interstitial Pneumonia}-- The histopathological pattern in IPF with temporal heterogeneity, fibroblast foci, and honeycombing. HRCT shows basal-predominant, subpleural reticular opacities with traction bronchiectasis.

\medterm{Ventilation-Perfusion Mismatch}-- The most common cause of hypoxemia. Areas of lung are ventilated but not perfused (dead space) or perfused but not ventilated (shunt). It is the mechanism of hypoxemia in COPD, pneumonia, and PE.

\medterm{Vocal Cord Dysfunction}-- Abnormal adduction of the vocal cords during inspiration, causing inspiratory stridor and dyspnea. It may mimic asthma but does not respond to bronchodilators. It is often misdiagnosed as refractory asthma. Diagnosis is by laryngoscopy showing paradoxical vocal cord motion. Treatment is speech therapy and avoiding triggers.

\medterm{Wheezes}-- High-pitched, continuous, musical sounds heard primarily during expiration, caused by airway narrowing. They are a hallmark of asthma and COPD. Expiratory wheezing indicates intrathoracic airway obstruction; inspiratory wheezing indicates extrathoracic obstruction. A silent chest (absent wheezing) in severe asthma indicates critical obstruction.

\medterm{Wheezing}-- A continuous, high-pitched, musical adventitious lung sound produced by turbulent airflow through narrowed or compressed lower airways during expiration (and occasionally inspiration). Pathophysiology: bronchospasm, mucosal edema, airway inflammation, or intraluminal secretions. Most common causes: Asthma, Chronic Obstructive Pulmonary Disease (COPD), Anaphylaxis, Heart Failure ("cardiac asthma"), and foreign body aspiration (focal monophonic wheeze).

\medterm{A-a Gradient}-- PAO2 minus PaO2, reflecting gas transfer efficiency. Normal approximately 15 mmHg (increases with age: < 15 + age/4). Elevated in V/Q mismatch, shunt, and diffusion impairment. Normal gradient in hypoventilation.

\medterm{Absent Breath Sounds}-- Complete absence of air movement sounds in a lung field, indicating complete airway obstruction, pneumothorax, or large pleural effusion. It is a medical finding requiring urgent evaluation and intervention.

\medterm{Alveolar Gas Equation}-- PAO2 = (FiO2 x (Patm - PH2O)) - (PaCO2 / R). Patm = 760 mmHg, PH2O = 47 mmHg, R = 0.8. Allows calculation of A-a gradient.

\medterm{Anaphylactic Shock}-- See immunology.

\medterm{Anoxia}-- Complete absence of oxygen supply to tissues, distinguished from hypoxia (reduced oxygen). Anoxia causes rapid cell death, particularly in highly aerobic tissues (brain, heart, kidneys). Causes: airway obstruction (complete—foreign body, laryngeal edema, anaphylaxis), respiratory arrest, severe drowning, carbon monoxide poisoning (carboxyhaemoglobin >50\%), and high-altitude exposure. Irreversible brain damage occurs within 4--6 minutes of complete anoxia (cerebral cortex most vulnerable, followed by hippocampus, basal ganglia, and cerebellum). EEG may show isoelectric activity. Hypoxic-ischemic encephalopathy (HIE) is the clinical syndrome following cardiac arrest or perinatal asphyxia. Treatment: immediate restoration of oxygenation (ABC, oxygen, ventilation), therapeutic hypothermia for HIE (reduces neurological injury), and supportive care.

\medterm{Apnea-Hypopnea Index}-- Number of apneas plus hypopneas per hour of sleep. Apnea = airflow cessation > 10 seconds. Hypopnea = partial obstruction with desaturation or arousal. AHI < 5 normal, 5--14 mild, 15--29 moderate, >= 30 severe.

\medterm{Apneustic Breathing}-- A breathing pattern characterized by a prolonged inspiratory pause, seen in brainstem lesions (particularly the pons). It indicates damage to the apneustic center and is a poor prognostic sign in neurological injury. It may require mechanical ventilation.

\medterm{Apnoea}-- Temporary cessation of breathing, defined as complete cessation of airflow for \(\ge\)10 seconds. Obstructive apnea: airway collapse despite respiratory effort (common in obstructive sleep apnea). Central apnea: absence of respiratory effort (Cheyne–Stokes respiration in heart failure, high-altitude periodic breathing, opioid-induced central apnea, brainstem stroke). Mixed apnea: features of both. Physiologic apnea: breath-holding, voluntary apnea. Pathological apnea in infants: apnea of prematurity (seen in preterm infants <37 weeks, due to immature respiratory control, treated with caffeine), apparent life-threatening events (ALTEs, now termed brief resolved unexplained events—BRUE). Diagnosis: polysomnography (sleep study). See also \hyperlink{Sleep Apnea}{Sleep apnea}.

\medterm{Aspiration Pneumonia}-- Lung infection from aspiration of oropharyngeal or gastric contents. Risk factors include impaired consciousness and dysphagia. Right lower lobe is most commonly affected.

\medterm{Assist-Control Ventilation}-- Ventilator delivers a set tidal volume at a set rate. Every breath receives the full set volume. Guarantees minimum rate and minute ventilation. Risk of respiratory alkalosis and auto-PEEP.

\medterm{Ataxic Breathing}-- An irregular, unpredictable breathing pattern with random variations in rate and depth, seen in medullary lesions. It is the most dangerous respiratory pattern and often precedes respiratory arrest. It requires mechanical ventilation.

\medterm{Atelectasis}-- Complete or partial collapse of lung tissue, resulting in loss of aeration. It may be caused by airway obstruction (resorption atelectasis), pleural disease (compressive atelectasis), surfactant dysfunction (dissolution atelectasis), or scarring (cicatrization atelectasis). It impairs gas exchange and may predispose to infection.

\medterm{Autogenic Drainage}-- A breathing technique using controlled breathing at different lung volumes to mobilize secretions from peripheral to central airways. It involves breathing at low lung volumes (collect), mid lung volumes (loosen), and high lung volumes (evacuate). It requires patient training and cooperation.

\medterm{Azygos Vein}-- A vein that drains the posterior intercostal veins and empties into the superior vena cava. It forms an important collateral pathway when the SVC is obstructed. On chest imaging, enlargement of the azygos vein may indicate SVC obstruction, heart failure, or portal hypertension.

\medterm{Barotrauma}-- Lung injury from excessive pressure during ventilation, causing alveolar rupture (pneumothorax, pneumomediastinum). Risk factors: high pressures, high PEEP. Prevented by lung-protective ventilation.

\medterm{Biot Breathing}-- An irregular breathing pattern with random periods of normal breathing interspersed with apneic episodes, seen in brainstem lesions. It is less organized than Cheyne-Stokes breathing and indicates medullary damage.

\medterm{Biotrauma}-- Inflammatory response to ventilator-induced lung injury with systemic cytokine release. Can lead to multi-organ dysfunction. Minimized by lung-protective ventilation strategies.

\medterm{Body Plethysmography}-- Technique for measuring lung volumes (TLC, RV, FRC) using a sealed chamber. Gold standard for TLC measurement. Useful for diagnosing restriction and hyperinflation.

\medterm{Breath Sounds}-- The sounds produced by airflow through the respiratory tract during breathing, assessed by auscultation with a stethoscope over the chest wall. Normal breath sounds are classified by location and quality: tracheal (loud, harsh, high-pitched heard over the trachea), bronchial (loud, tubular, heard over the manubrium and large airways, with a gap between inspiration and expiration), bronchovesicular (intermediate pitch and intensity, heard over the mainstem bronchi, equal inspiratory and expiratory duration,.

\medterm{Breath Stacking}-- A technique using a manual resuscitation bag or glossopharyngeal breathing to deliver successive breaths without exhaling, stacking tidal volumes to achieve higher lung volumes. It improves cough effectiveness and secretion clearance in patients with weak respiratory muscles.

\medterm{Bronchial Stenosis}-- Narrowing of a bronchus from trauma, infection, inflammation, or prior intervention. It causes localized wheezing, recurrent infections, and impaired secretion clearance. Diagnosis is by bronchoscopy. Treatment includes balloon dilation, stenting, or surgical resection.

\medterm{Bronchogenic Cyst}-- A congenital cyst derived from abnormal budding of the ventral foregut. It is typically found in the middle mediastinum or near the carina. It is usually asymptomatic but may cause compression of adjacent structures. Treatment is surgical resection if symptomatic.

\medterm{Bronchophony}-- Increased transmission of spoken sounds to the stethoscope over consolidated lung tissue. It is normally heard over the central airways. Excessive bronchophony over peripheral lung fields indicates consolidation. It is part of the constellation of physical examination findings in pneumonia.

\medterm{Bronchopneumonia}-- Patchy consolidation centered on bronchioles across multiple lobes. Caused by S. aureus, gram-negative bacteria, and Pseudomonas.

\medterm{Bronchoscopy}-- See medical\_tests.

\medterm{Carina}-- The ridge at the base of the trachea where it bifurcates into the right and left main bronchi. It is a landmark for endotracheal tube placement and bronchoscopy. The right main bronchus is wider, shorter, and more vertical than the left, making it the preferred path for aspirated foreign bodies.

\medterm{Central Sleep Apnea}-- Apneas from variability in respiratory effort without airway obstruction. Associated with heart failure (Cheyne-Stokes respiration), high altitude, and opioid use.

\medterm{CFTR Modulator Therapy}-- Targeted therapies correcting the defective CFTR protein. Elexacaftor-tezacaftor-ivacaftor (Trikafta) is the most effective combination for patients with at least one F508del mutation. These therapies have dramatically changed CF prognosis.

\medterm{Chest Tube Insertion}-- Placement of a thoracostomy tube in the safe triangle (latissimus dorsi, pectoralis major, nipple level) to drain air, fluid, or blood from the pleural space. Connected to underwater seal drainage.

\medterm{Chest Wall Deformity}-- Abnormalities of the chest wall that can affect lung function. Pectus excavatum (sunken sternum) may compress the heart and lungs. Pectus carinatum (protruding sternum) is usually cosmetic. Scoliosis can cause restrictive lung disease. These may require surgical correction if severe.

\medterm{Cheyne-Stokes Respiration}-- Crescendo-decrescendo breathing pattern with periods of apnea, seen in heart failure and stroke. It is a form of central sleep apnea.

\medterm{Chlamydophila Pneumonia}-- Atypical pneumonia with mild presentation: sore throat, hoarseness, cough. Treatment is macrolides or tetracyclines.

\medterm{Chylothorax}-- Pleural effusion with chyle from thoracic duct disruption. Triglycerides > 110 mg/dL. Treatment includes dietary modification, octreotide, and surgical ligation if refractory.

\medterm{Community-Acquired Pneumonia}-- See infectious\_disease.

\medterm{Compressive Atelectasis}-- Atelectasis caused by external compression of the lung by pleural fluid, air, or a mass. The lung volume is reduced by the external force. It is reversible when the compressing force is relieved (e.g., drainage of pleural effusion).

\medterm{Cor Pulmonale}-- Right ventricular hypertrophy and dilation from pulmonary hypertension due to lung disease. Features include peripheral edema, JVD, and hepatomegaly.

\medterm{Cough}-- A complex reflex arc protecting the airway by clearing secretions, irritants, and foreign material from the tracheobronchial tree. The cough reflex involves sensory receptors (rapidly adapting irritant receptors, C-fibers, and mechanoreceptors) in the larynx, trachea, and bronchi, afferent pathways (vagus nerve, superior laryngeal nerve, trigeminal nerve), the cough center in the medulla oblongata, efferent pathways (recurrent laryngeal nerve, phrenic nerve, spinal motor nerves), and effector muscles (diaphragm, abdominal wall, intercostals, and laryngeal muscles) for the coordinated cough response.

\medterm{Cough Augmentation}-- Techniques to improve cough effectiveness in patients with neuromuscular disease or chest wall weakness. Methods include manually assisted cough (abdominal thrust timed with cough), mechanical insufflation-exsufflation (CoughAssist device), and breath stacking with a resuscitation bag. They are essential in Duchenne muscular dystrophy and spinal cord injury.

\medterm{CPAP}-- Continuous Positive Airway Pressure, first-line for OSA. Delivers constant pressure through a mask to prevent airway collapse. Eliminates apneas and improves outcomes. Adherence is the main challenge.

\medterm{Crazy Paving Pattern}-- A HRCT finding characterized by ground-glass opacification with superimposed interlobular septal thickening and intralobular lines. It is seen in pulmonary alveolar proteinosis, lipoid pneumonia, ARDS, pulmonary hemorrhage, and Pneumocystis pneumonia. Despite its name, it is not specific to any single diagnosis.

\medterm{Cryptogenic Organizing Pneumonia}-- An inflammatory lung disease with granulation tissue within alveolar ducts. May be idiopathic or secondary to infection, drugs, or connective tissue disease. Responds well to corticosteroids.

\medterm{CT Pulmonary Angiography}-- Gold standard for PE diagnosis, showing thrombi as filling defects. Sensitivity > 95\%, specificity > 97\% for central emboli.

\medterm{CURB-65 Score}-- One point each for: Confusion, Urea > 7 mmol/L, RR > 30, BP < 90/60, age >= 65. 0--1: outpatient. 2: consider hospitalization. 3--5: severe.

\medterm{D-Dimer}-- See medical\_tests.

\medterm{Decreased Breath Sounds}-- Diminished air movement heard on auscultation, indicating decreased ventilation. Causes include pneumothorax, pleural effusion, atelectasis, and mucus plugging. Unilateral decreased breath sounds may indicate tension pneumothorax or large effusion.

\medterm{Directly Observed Therapy}-- Supervision of medication ingestion to ensure adherence. Standard of care for TB treatment. Ensures completion of the full treatment course and reduces development of drug resistance.

\medterm{Dornase Alfa}-- A recombinant human DNase that cleaves extracellular DNA in airway secretions, reducing sputum viscosity. It is used in CF to improve lung function and reduce exacerbations. It is administered as a nebulized solution once or twice daily. It has become a standard part of CF airway clearance regimens.

\medterm{Dry Powder Inhaler}-- A breath-actuated inhaler delivering powdered medication. Does not require hand-breath coordination. Requires adequate inspiratory flow. Examples include Advair Diskus and Spiriva HandiHaler.

\medterm{Egophony}-- An abnormal vocal sound heard during auscultation in which spoken ``E'' sounds are heard as ``A'' (the ``E-to-A'' change). It is a sign of lung consolidation, as in lobar pneumonia. The consolidated lung tissue conducts high-frequency sounds better than air-filled lung.

\medterm{Electromagnetic Navigation Bronchoscopy}-- A guidance system for bronchoscopy that uses electromagnetic fields to navigate to peripheral lung lesions. It allows biopsy of lesions not visible on standard bronchoscopy. It is more accurate than conventional transbronchial biopsy for peripheral lesions. It may reduce the need for CT-guided biopsy.

\medterm{Empyema}-- Pus in the pleural space. CT shows peripheral enhancement (split pleura sign). Treatment requires chest tube drainage, intrapleural fibrinolytics, or VATS decortication.

\medterm{Empyema Necessitans}-- Extension of pleural empyema through the chest wall. Presents as a painful, fluctuant mass. Requires drainage and antibiotics.

\medterm{Endobronchial Tumor}-- A tumor growing within the bronchial lumen, causing obstruction, atelectasis, and post-obstructive pneumonia. It may be malignant (lung cancer, carcinoid tumor) or benign (hamartoma, papilloma). Symptoms include cough, hemoptysis, and wheezing. Treatment is bronchoscopic removal or debulking plus definitive therapy.

\medterm{Endobronchial Ultrasound}-- Bronchoscopy with ultrasound guidance for transbronchial needle aspiration of mediastinal and hilar lymph nodes. Essential for lung cancer staging. Less invasive than mediastinoscopy.

\medterm{Endothelin Receptor Antagonists}-- Medications that block endothelin-1, a potent vasoconstrictor implicated in pulmonary arterial hypertension. Bosentan (dual ERA), ambrisentan (selective ETA), and macitentan are used for PAH. Side effects include hepatotoxicity (bosentan), anemia, and teratogenicity.

\medterm{Endotracheal Intubation}-- Placement of a tube through the mouth into the trachea for mechanical ventilation. Gold standard for airway management. Confirmation by end-tidal CO2 detection and bilateral breath sounds.

\medterm{Expiratory Grunting}-- A forced exhalation against a partially closed glottis, creating an audible sound. It is an end-expiratory maneuver that generates intrinsic PEEP, helping to maintain alveolar recruitment and improve oxygenation. It is commonly seen in children with respiratory distress.

\medterm{Extensively Drug-Resistant TB}-- TB resistant to isoniazid, rifampin, fluoroquinolones, and at least one injectable agent. Very limited treatment options and poor outcomes.

\medterm{Exudative Pleural Effusion}-- From local pleural disease. Causes include pneumonia, malignancy, TB, and PE. Elevated protein, LDH, and cell count.

\medterm{Farmer's Lung}-- The most common HP, caused by thermophilic actinomycetes in moldy hay. Presents acutely 4--8 hours after exposure with fever, cough, and dyspnea.

\medterm{Flail Chest}-- A clinical condition caused by three or more consecutive ribs fractured in two or more places, creating a free-floating segment of chest wall. The segment moves paradoxically inward during inspiration and outward during expiration. It causes severe pain, impaired ventilation, and may require mechanical ventilation. It is often associated with pulmonary contusion.

\medterm{Flow-Volume Loop}-- Graphical representation of airflow during forced breathing. Obstructive: scooped expiratory limb. Restrictive: narrowed loop with preserved shape. Variable extrathoracic obstruction: flattening of inspiratory limb.

\medterm{Foreign Body Aspiration}-- See emergency\_medicine.

\medterm{Hamman's Sign}-- A crunching or crackling sound heard on auscultation synchronized with the heartbeat, caused by air in the mediastinum (pneumomediastinum). It is a rare but specific finding. It is also described as a mediastinal crunch heard in spontaneous pneumomediastinum.

\medterm{Hampton Hump}-- Wedge-shaped opacification with its base abutting the pleura on chest radiography, representing pulmonary infarction secondary to Pulmonary Embolism. \textit{See Pulmonary Embolism (PE).}

\medterm{Hampton's Hump}-- Wedge-shaped, pleural-based opacity on CXR representing pulmonary infarction distal to PE.

\medterm{Hemothorax}-- Blood in pleural space from trauma. Initial chest tube output > 1500 mL or persistent bleeding > 200 mL/hour indicates need for surgery.

\medterm{Hereditary Hemorrhagic Telangiectasia}-- An autosomal dominant vascular disorder (also called Osler-Weber-Rendu syndrome) characterized by telangiectasias on skin and mucous membranes, pulmonary AVMs, hepatic AVMs, and cerebral AVMs. Pulmonary AVMs can cause hypoxemia and paradoxical emboli. Treatment is embolization of pulmonary AVMs.

\medterm{High-Flow Nasal Cannula}-- A device delivering heated and humidified oxygen at flow rates up to 60 L/min through a nasal cannula. It provides a modest amount of PEEP, reduces dead space ventilation, and improves patient comfort compared to non-invasive ventilation. It is used for acute hypoxemic respiratory failure and post-extubation support.

\medterm{Hospital-Acquired Pneumonia}-- Pneumonia developing > 48 hours after hospital admission. Common pathogens include Pseudomonas and MRSA. Empiric treatment covers resistant organisms.

\medterm{Hypercyanotic Tetralogy}-- A complex congenital heart defect with four features: ventricular septal defect, overriding aorta, right ventricular outflow tract obstruction (pulmonary stenosis), and right ventricular hypertrophy. It causes right-to-left shunting and cyanosis. Treatment is surgical repair, typically in infancy.

\medterm{Hyperventilation}-- Breathing in excess of metabolic requirements, resulting in decreased partial pressure of arterial carbon dioxide (PaCO2) and respiratory alkalosis. Acute hyperventilation causes: perioral and digital paraesthesiae (numbness, tingling of lips, hands, feet), carpopedal spasm (tetany due to decreased ionised calcium from alkalosis--Chvostek and Trousseau signs), lightheadedness, dizziness, chest tightness, palpitations, anxiety, panic, and rarely loss of consciousness. Causes: anxiety and panic disorder (most common, acute or chronic hyperventilation syndrome), pain, fever, sepsis, pulmonary embolism (tachypnoea, hypoxia, hyperventilation), head injury (lesions affecting the respiratory center), metabolic acidosis (Kussmaul breathing as compensatory hyperventilation, e.g., DKA), drugs (salicylate overdose), pregnancy (progesterone stimulates the respiratory center), and neurological disorders (stroke, brainstem lesions). Diagnosis: clinical (typical symptoms, normal oxygen saturation, history of anxiety/panic), arterial blood gas (decreased PaCO2 <35 mmHg, elevated pH >7.45, normal PO2). Hyperventilation can be distinguished from pathological tachypnoea by the absence of hypoxia and underlying cause. Management: (1) Acute: reassurance, rebreathing into a paper bag (controversial and risk of hypoxia safer to use slow breathing technique), breathing retraining (5-second inhale, 5-second hold, 5-second exhale), treatment of underlying cause. (2) Chronic: cognitive behavioural therapy (CBT), breathing exercises, management of anxiety/panic disorder (SSRI, buspirone, relaxation techniques).

\medterm{Hypoxemia}-- PaO2 < 80 mmHg on room air. Causes: hypoventilation, V/Q mismatch, shunt, diffusion impairment, low FiO2. A-a gradient helps distinguish causes.

\medterm{Influenza}-- A contagious respiratory illness caused by influenza A or B viruses, characterized by sudden onset of fever, myalgia, headache, malaise, nonproductive cough, sore throat, and nasal congestion. Seasonal epidemics occur annually. Complications: pneumonia (primary viral or secondary bacterial), exacerbation of underlying conditions. Diagnosis: PCR or rapid antigen testing. Treatment: supportive care, neuraminidase inhibitors (oseltamivir). Prevention: annual vaccination.

\medterm{Inhaled Nitric Oxide}-- A selective pulmonary vasodilator used in acute settings (acute respiratory failure, right heart failure, post-cardiac surgery). It reduces pulmonary vascular resistance and improves oxygenation by dilating well-ventilated pulmonary vessels (improving V/Q matching). It has a short half-life and requires continuous delivery.

\medterm{Interferon-Gamma Release Assay}-- Blood tests measuring IFN-gamma release in response to M. tuberculosis antigens (ESAT-6, CFP-10). More specific than PPD and unaffected by BCG. Cannot distinguish latent from active TB.

\medterm{Intrapulmonary Shunt}-- Perfusion of alveoli without ventilation (V/Q = 0). Blood passes through the pulmonary circulation without gas exchange, mixing with oxygenated blood and causing hypoxemia. It is seen in ARDS, pneumonia, and atelectasis. It is refractory to supplemental oxygen.

\medterm{IVC Filter}-- A device in the IVC to prevent DVT from reaching the lungs. Indications include recurrent PE despite anticoagulation and absolute contraindication to anticoagulation.

\medterm{Klebsiella Pneumonia}-- Severe, necrotizing pneumonia in alcoholic or debilitated patients. Upper lobe involvement with ``bulging fissure sign'' and ``currant jelly'' sputum.

\medterm{Kussmaul Breathing}-- A deep, rapid breathing pattern seen in severe metabolic acidosis (diabetic ketoacidosis, lactic acidosis, uremia). It represents the body's attempt to compensate by blowing off CO2. The increased respiratory rate and depth create the characteristic labored breathing.

\medterm{Kyphoscoliosis}-- An abnormal curvature of the spine in both the coronal (scoliosis) and sagittal (kyphosis) planes. Severe forms cause restrictive lung disease by reducing chest wall compliance and lung volumes. It can lead to chronic respiratory failure and pulmonary hypertension. Treatment includes bracing, surgical correction, and non-invasive ventilation.

\medterm{Lambert-Eaton Myasthenic Syndrome}-- Autoantibodies against presynaptic calcium channels at the neuromuscular junction, associated with SCLC. Proximal weakness, reduced reflexes, facilitation with exercise.

\medterm{Large Airway Obstruction}-- Obstruction of airways with diameter > 2 mm (trachea and main bronchi). It causes significant symptoms including stridor, dyspnea, and impaired secretion clearance. It may be fixed (stenosis, tumor) or variable (tracheomalacia, vocal cord dysfunction). Flow-volume loops show characteristic patterns.

\medterm{Latent Tuberculosis Infection}-- Persistent immune response to M. tuberculosis without active disease. Asymptomatic and non-infectious. 5--10\% of immunocompetent individuals develop active TB without treatment. Treatment reduces risk by 60--90\%.

\medterm{Legionella Pneumonia}-- From Legionella pneumophila in water systems. Presents with high fever, diarrhea, confusion, and hyponatremia. Diagnosis by urinary antigen test. Treatment is fluoroquinolones or macrolides.

\medterm{Light's Criteria}-- An effusion is exudative if any one is met: pleural protein/serum protein > 0.5, pleural LDH/serum LDH > 0.6, or pleural LDH > 2/3 upper limit of normal serum LDH.

\medterm{Lipoid Pneumonia}-- Lung inflammation caused by accumulation of lipid-laden macrophages in the alveoli. Exogenous lipoid pneumonia results from aspiration of mineral oil, petroleum jelly, or oil-based nose drops. Endogenous lipoid pneumonia occurs distal to bronchial obstruction. HRCT shows ground-glass opacities with low-attenuation areas (fat density).

\medterm{Lobar Pneumonia}-- Consolidation involving an entire lobe. Most commonly caused by S. pneumoniae. Classic progression: congestion, red hepatization, gray hepatization, resolution.

\medterm{Low-Dose CT Scan}-- CT protocol using reduced radiation for lung cancer screening. Detects small nodules at earlier stages. False positives require follow-up per Fleischner guidelines.

\medterm{Lung Abscess}-- Localized necrosis with cavitation containing purulent material. Causes include aspiration and anaerobic bacteria. CXR shows thick-walled cavity with air-fluid level.

\medterm{Lung Function Tests}-- A group of physiological tests that assess the function and capacity of the lungs. See \hyperlink{Pulmonary Function Tests}{Pulmonary Function Tests}.

\medterm{Lung Protective Ventilation}-- A ventilation strategy using low tidal volumes (6 mL/kg ideal body weight), plateau pressure < 30 cmH2O, appropriate PEEP, and permissive hypercapnia to minimize ventilator-induced lung injury. It is the standard of care for ARDS and is used in all mechanically ventilated patients to reduce barotrauma, volutrauma, and biotrauma.

\medterm{Lymphangioleiomyomatosis}-- A rare, progressive cystic lung disease caused by proliferation of abnormal smooth muscle-like cells (LAM cells) in the lungs, lymphatics, and kidneys. It occurs almost exclusively in women of childbearing age. HRCT shows thin-walled cysts throughout both lungs. It may be associated with tuberous sclerosis complex. Treatment includes sirolimus (mTOR inhibitor) and lung transplantation.

\medterm{Massive Hemoptysis}-- > 600 mL in 24 hours or causing hemodynamic compromise. Medical emergency. Treatment: airway protection, bronchial artery embolization (first-line), or surgery.

\medterm{Mechanical Insufflation-Exsufflation}-- A device that delivers positive pressure (insufflation) followed by rapid negative pressure (exsufflation) to simulate a cough. It is used in neuromuscular disease (Duchenne muscular dystrophy, spinal cord injury, ALS) to augment cough effectiveness and clear secretions.

\medterm{Mediastinoscopy}-- Surgical procedure to sample mediastinal lymph nodes through a small incision above the sternum. Gold standard for mediastinal lymph node staging in lung cancer. Being replaced by EBUS in many centers.

\medterm{Miliary TB}-- Hematogenous dissemination of M. tuberculosis causing innumerable 1--3 mm granulomas throughout lungs and other organs. CXR shows ``millet seed'' pattern. More common in immunocompromised patients.

\medterm{Minute Ventilation}-- The total volume of air breathed in one minute (tidal volume x respiratory rate). Normal is 6--8 L/min. Elevated in increased metabolic demand, pain, anxiety, and sepsis. Reduced in neuromuscular disease and sedation.

\medterm{Multidrug-Resistant TB}-- TB resistant to at least isoniazid and rifampin. Requires second-line drugs for 18--24 months. Worse outcomes and higher mortality than drug-susceptible TB.

\medterm{Mycoplasma Pneumonia}-- Atypical pneumonia most common in ages 5--35. Dry cough, low-grade fever, headache. Patchy infiltrates on CXR disproportionate to symptoms (``walking pneumonia'').

\medterm{Negative Pressure Ventilation}-- Ventilation using negative pressure applied around the chest wall to expand the lungs, as opposed to positive pressure applied to the airway. The ``iron lung'' is a historical example. Modern negative pressure ventilators include the cuirass and wrap. They are rarely used today but may be useful in neuromuscular disease.

\medterm{Obesity Hypoventilation Syndrome}-- Chronic daytime hypercapnia in obese individuals without other causes. Reduced chest wall compliance and blunted respiratory drive. Often coexists with OSA. Treatment: weight loss and BiPAP.

\medterm{Obstructive Sleep Apnea}-- Repetitive upper airway obstruction during sleep. Presents with snoring, witnessed apneas, daytime sleepiness. AHI quantifies severity: 5--14 mild, 15--29 moderate, >= 30 severe. Treatment: CPAP, oral appliances, weight loss, surgery.

\medterm{Orthopnea}-- Dyspnea that occurs when lying flat and is relieved by sitting upright. It is measured by the number of pillows the patient needs to sleep (two-pillow orthopnea). It is a hallmark of left heart failure and is caused by redistribution of fluid from the lower extremities to the lungs in the supine position.

\medterm{Pancoast Tumor}-- Superior sulcus tumor invading brachial plexus, subclavian vessels, and sympathetic chain. Causes Horner syndrome, shoulder/arm pain, and hand muscle atrophy. Treated with chemoradiation followed by surgery.

\medterm{Parapneumonic Effusion}-- Pleural effusion complicating pneumonia. Simple effusions resolve with antibiotics. Complicated effusions (pH < 7.2, glucose < 60 mg/dL) require chest tube drainage.

\medterm{Paroxysmal Nocturnal Dyspnea}-- Episodes of severe dyspnea that awaken the patient from sleep, typically occurring 1--2 hours after lying down. It is a classic symptom of left heart failure, caused by progressive pulmonary congestion during sleep. It is different from orthopnea in its delayed onset and nocturnal presentation.

\medterm{PE Anticoagulation}-- Initial heparin (LMWH or UFH), followed by warfarin (INR 2--3) for 3--6 months. DOACs are first-line for most patients. Duration depends on provoked vs. unprovoked status.

\medterm{Permissive Hypercapnia}-- A ventilator strategy that allows PaCO2 to rise above normal (and pH to fall) in exchange for lung-protective ventilation with lower tidal volumes and airway pressures. It is used in severe ARDS when achieving normal PaCO2 would require harmful ventilation parameters.

\medterm{Phosphodiesterase-5 Inhibitors}-- Medications that inhibit PDE-5, increasing cGMP and causing smooth muscle relaxation and vasodilation. Sildenafil and tadalafil are used for PAH and erectile dysfunction. In PAH, they reduce pulmonary vascular resistance and improve exercise capacity.

\medterm{Platypnea}-- Dyspnea that occurs in the upright position and is relieved by lying down. It is the opposite of orthopnea. It is rare and may be associated with hepatopulmonary syndrome, intracardiac shunting, or pulmonary arteriovenous malformations.

\medterm{Pleural Effusion}-- Abnormal fluid in the pleural space. Transudates result from altered pressures; exudates from inflammation, infection, or malignancy. Light's criteria differentiate the two.

\medterm{Pleural Fluid Glucose}-- Low pleural fluid glucose (< 60 mg/dL) suggests bacterial empyema, rheumatoid pleurisy, or tuberculosis. It indicates high metabolic activity in the pleural space consuming glucose. Very low glucose is characteristic of empyema and parapneumonic effusions requiring drainage.

\medterm{Pleural Fluid pH}-- Pleural fluid pH < 7.20 indicates a complicated parapneumonic effusion or empyema requiring drainage. Normal pleural fluid pH is approximately 7.60. Low pH reflects high metabolic activity and anaerobic metabolism in the pleural space. It is a key decision point in managing parapneumonic effusions.

\medterm{Pleural Friction Rub}-- A grating, leathery sound heard on auscultation caused by inflamed pleural surfaces rubbing against each other during breathing. It is a sign of pleurisy and may be heard in pneumonia, PE, and pleural effusion. It disappears when fluid accumulates and separates the pleural surfaces.

\medterm{Pleural Plaque}-- Circumscribed fibrous tissue on parietal pleura, hallmark of asbestos exposure. Bilateral, often calcified. Does not indicate asbestosis.

\medterm{Pleurisy}-- Pleural inflammation causing sharp, breathing-related chest pain. Causes: pneumonia, PE, autoimmune disease, malignancy. Treatment addresses the underlying cause.

\medterm{Pleuritis}-- Inflammation of the pleural membranes (visceral and parietal pleura). Presents with pleuritic chest pain: sharp, stabbing pain exacerbated by deep breathing, coughing, or sneezing. Common causes: viral pleurisy (most common), bacterial pneumonia with parapneumonic effusion, pulmonary embolism, tuberculosis, connective tissue diseases (SLE, rheumatoid arthritis), uremia, and malignancy. May be accompanied by pleural friction rub on auscultation. Diagnosis: clinical history, chest X-ray (evaluate for effusion, infiltrates, or pneumothorax), ultrasound, and laboratory evaluation (CBC, CRP, D-dimer to rule out PE). Treatment: addressing the underlying cause; NSAIDs for pain relief; corticosteroids for severe inflammatory pleuritis (e.g., SLE).

\medterm{Pleurodesis}-- Fusion of visceral and parietal pleura to prevent recurrence of effusions or pneumothorax. Chemical pleurodesis uses talc, doxycycline, or bleomycin.

\medterm{Pneumatocele}-- Thin-walled, air-filled cystic spaces in the lung, often in staphylococcal pneumonia. Typically self-resolving.

\medterm{Pneumocystis Pneumonia}-- Opportunistic pneumonia from P. jirovecii in immunocompromised patients (CD4 < 200). Progressive dyspnea, dry cough, bilateral ground-glass opacities. Elevated LDH. Treatment is TMP-SMX.

\medterm{Pneumonia}-- Infection of the lung parenchyma, classified as community-acquired (CAP), hospital-acquired (HAP), or ventilator-associated (VAP). Common pathogens: Streptococcus pneumoniae, Haemophilus influenzae, Mycoplasma pneumoniae, viruses. Presents with cough, fever, dyspnea, and consolidation on imaging. Severity assessment: CURB-65 or PSI. Treatment: empiric antibiotics based on setting and local resistance patterns.

\medterm{Pneumothorax}-- Air in the pleural space causing lung collapse. Primary spontaneous occurs in tall, thin young men. Secondary occurs with underlying lung disease. Tension is a medical emergency.

\medterm{Positive End-Expiratory Pressure}-- Pressure maintained at end-expiration to prevent alveolar collapse. Improves oxygenation by recruiting alveoli and increasing FRC. Balances oxygenation with hemodynamic effects.

\medterm{Post-Infectious Cough}-- A chronic cough persisting 3--8 weeks after a URI, caused by transient bronchial hyperresponsiveness and airway inflammation. It is self-limited and usually resolves without treatment. Inhaled corticosteroids, antitussives, and time are the mainstays of management.

\medterm{Post-Obstructive Pneumonia}-- Pneumonia occurring distal to a bronchial obstruction caused by tumor, foreign body, mucus plug, or stenosis. The obstruction prevents clearance of secretions, creating a favorable environment for infection. It does not resolve with antibiotics alone until the obstruction is relieved.

\medterm{Post-Pneumonectomy Syndrome}-- A rare complication following pneumonectomy where the remaining lung herniates across the mediastinum, causing compression of airways and vessels. It presents with dyspnea, stridor, and recurrent infections. Treatment is surgical correction with mediastinal repositioning and prosthetic implant placement.

\medterm{Pott Disease}-- TB of the spine involving thoracic or lumbar vertebrae. Causes vertebral destruction, kyphosis (gibbus deformity), and paravertebral abscess. May cause spinal cord compression.

\medterm{Pressure Support Ventilation}-- Spontaneous breathing mode with ventilator-provided pressure support above PEEP for each breath. Reduces work of breathing. Used for weaning from mechanical ventilation.

\medterm{Primary Tuberculosis}-- Initial M. tuberculosis infection causing Ghon focus in lower lobes with hilar lymph node involvement (Ghon complex). Most infections are contained, resulting in latent TB.

\medterm{Prone Positioning}-- Turning supine patients to prone for > 12 hours daily in severe ARDS. Improves oxygenation by recruiting dorsal lung segments. Reduces mortality when PaO2/FiO2 < 150.

\medterm{Prostacyclin Analogs}-- Medications that mimic prostacyclin (PGI2), causing vasodilation and inhibiting platelet aggregation in the pulmonary vasculature. They include epoprostenol (continuous IV), treprostinil (subcutaneous, IV, inhaled, oral), and iloprost (inhaled). They are used for severe PAH.

\medterm{Pulmonary Alveolar Hemorrhage}-- Bleeding into the alveolar spaces from disruption of the alveolar-capillary membrane. Causes include vasculitis (GPA, MPA, Goodpasture), connective tissue disease, medications, and coagulopathy. It presents with hemoptysis, falling hemoglobin, diffuse bilateral infiltrates, and serially bloody returns on BAL.

\medterm{Pulmonary Alveolar Microlithiasis}-- A rare genetic disease caused by SLC34A2 mutations, leading to accumulation of calcium phosphate microliths within the alveoli. HRCT shows diffuse bilateral calcified nodules with a ``sandstorm'' appearance. It is usually diagnosed in the third to fourth decade. Treatment is supportive; lung transplantation may be needed for progressive disease.

\medterm{Pulmonary Alveolar Proteinosis}-- A rare disease characterized by accumulation of surfactant-like material within the alveoli due to impaired clearance by alveolar macrophages. Most cases are autoimmune (anti-GM-CSF antibodies). Presents with progressive dyspnea and cough. HRCT shows bilateral ground-glass opacities with ``crazy paving'' pattern. Treatment is whole lung lavage.

\medterm{Pulmonary Arteriovenous Malformation}-- An abnormal direct communication between a pulmonary artery and vein, causing a right-to-left shunt. It may be idiopathic or associated with hereditary hemorrhagic telangiectasia (Osler-Weber-Rendu syndrome). It can cause hypoxemia, paradoxical emboli (stroke, brain abscess), and hemoptysis. Treatment is embolization or surgical resection.

\medterm{Pulmonary Artery Catheter}-- A catheter inserted through the right heart into the pulmonary artery to measure hemodynamic parameters (PA pressure, PAWP, cardiac output). Used for diagnosis and management of pulmonary hypertension, heart failure, and critical illness.

\medterm{Pulmonary Contusion}-- Bruising of the lung parenchyma from blunt chest trauma, causing hemorrhage and edema without lung laceration. It impairs gas exchange and can progress to ARDS. It is often associated with rib fractures and flail chest. Treatment is supportive with lung-protective ventilation.

\medterm{Pulmonary Edema}-- Accumulation of fluid in the pulmonary interstitium and alveoli, most commonly due to left-sided heart failure (cardiogenic). Causes increased hydrostatic pressure leads to transudation. Presents with acute dyspnea, orthopnea, crackles, and S3 gallop. Non-cardiogenic causes include ARDS, renal failure, and high altitude. Treatment: diuretics, afterload reduction, oxygen, and positive pressure ventilation if severe.

\medterm{Pulmonary Embolism}-- See emergency\_medicine.

\medterm{Pulmonary Hemosiderosis}-- Iron deposition in the lungs from repeated alveolar hemorrhage. It presents with hemoptysis, anemia, and diffuse bilateral infiltrates on imaging. Causes include Goodpasture syndrome, granulomatosis with polyangiitis, and idiopathic pulmonary hemosiderosis. The classic triad is hemoptysis, anemia, and diffuse pulmonary infiltrates.

\medterm{Pulmonary Infarction}-- Lung necrosis from PE when bronchial supply is compromised. Causes pleuritic pain, hemoptysis, and wedge-shaped opacity. Most PE do not cause infarction due to dual blood supply.

\medterm{Pulmonary Sequestration}-- A congenital malformation of non-functioning lung tissue that lacks normal communication with the tracheobronchial tree and receives its blood supply from a systemic artery (usually the descending aorta). Intralobar sequestration is within the normal pleural investment; extralobar is outside. It may present with recurrent infections or be incidental.

\medterm{Pulmonary Vasodilator Therapy}-- Medications that reduce pulmonary vascular resistance and pulmonary artery pressure. Used for WHO Group 1 PAH. Includes endothelin receptor antagonists (bosentan), PDE-5 inhibitors (sildenafil), prostacyclin analogs (epoprostenol), and soluble guanylate cyclase stimulators (riociguat).

\medterm{Pulmonary Wedge Pressure}-- The pressure measured by a balloon-tipped pulmonary artery catheter when the balloon is inflated, occluding a branch of the pulmonary artery. It approximates left atrial pressure. Normal is 6--12 mmHg. Elevated in left heart failure (PCWP > 18 mmHg). Low in hypovolemia.

\medterm{Reactivation Tuberculosis}-- Reactivation of latent infection, typically years later. Characteristically affects apical and posterior upper lobe segments. Features include chronic cough, hemoptysis, night sweats, weight loss, and cavitation.

\medterm{Relapsing Polychondritis}-- A systemic autoimmune disorder causing inflammation and destruction of cartilage throughout the body. Airway involvement can cause tracheal and bronchial stenosis. Other features include auricular chondritis, nasal chondritis, ocular inflammation, and arthritis. Treatment is corticosteroids and immunosuppressants.

\medterm{Resorption Atelectasis}-- Atelectasis caused by complete obstruction of an airway, leading to absorption of trapped air by the pulmonary capillaries. Common causes include mucus plugging, foreign body, tumor, and post-surgical shallow breathing. The affected area appears as increased density on CXR with volume loss.

\medterm{RIPE Therapy}-- Standard 4-drug regimen for active TB: Rifampin, Isoniazid, Pyrazinamide, Ethambutol. Intensive phase 2 months, continuation phase 4 months (total 6 months). Pyridoxine prevents isoniazid-induced neuropathy.

\medterm{Round Pneumonia}-- Pneumonia appearing as a round opacity on CXR rather than typical consolidation. More common in children and elderly. May mimic lung mass.

\medterm{Siderosis}-- Benign pneumoconiosis from iron dust, seen in welders. Causes small opacities on CXR but minimal impairment. Reversible with removal of exposure.

\medterm{Situs Inversus}-- A congenital condition in which the thoracic and abdominal organs are reversed or mirrored from their normal positions. Complete situs inversus (situs inversus totalis) involves all organs. It is usually asymptomatic but may be associated with primary ciliary dyskinesia (Kartagener syndrome), congenital heart disease, and polysplenia or asplenia.

\medterm{Sleep Apnea (Obstructive Sleep Apnea - OSA)}-- A sleep disorder characterized by repeated episodes of partial or complete upper airway obstruction during sleep (obstructive sleep apnea—OSA, most common) or cessation of respiratory drive (central sleep apnea—CSA), causing intermittent hypoxia, hypercapnia, and sleep fragmentation. OSA affects ~13\% of men and 6\% of women (moderate--severe). Risk factors: obesity (BMI \(\ge\)30—strongest risk factor, 60--90\% of OSA patients), male sex, age >40, menopause, craniofacial abnormalities (retrognathia, macroglossia, tonsillar hypertrophy), smoking, alcohol, and hypothyroidism. Pathophysiology: loss of pharyngeal muscle tone during sleep leads to collapse of the pharyngeal airway (velopharynx, oropharynx, hypopharynx). Clinical: loud snoring, witnessed apnoeas (by bed partner), choking/gasping arousals, unrefreshing sleep, excessive daytime sleepiness (Epworth Sleepiness Scale \(\ge\)10), morning headache, nocturia, and cognitive impairment. Consequences: hypertension (dose-dependent, 40\% of OSA patients have hypertension), cardiovascular disease (MI, stroke, atrial fibrillation, heart failure), type 2 diabetes, and increased all-cause mortality. Diagnosis: polysomnography (apnea-hypopnoea index—AHI: mild 5--15, moderate 15--30, severe >30 events/hour). Home sleep apnea testing (HSAT) is an alternative for high-risk patients. Treatment: continuous positive airway pressure (CPAP—first-line, splints the airway open, AHI reduction >90\%), mandibular advancement devices (MAD—for mild--moderate OSA, <25\% of CPAP efficacy), weight loss (5--10\% weight loss reduces AHI by 30--50\%), positional therapy (for supine-dependent OSA), and surgery (uvulopalatopharyngoplasty, maxillomandibular advancement, hypoglossal nerve stimulation/Inspire—for moderate--severe OSA who fail CPAP). Untreated severe OSA: 5-year mortality ~10--15\%.

\medterm{Snoring and Sleep Apnea}-- Snoring is the sound produced by vibrating soft tissues in the upper airway during sleep, often resulting from partial airway obstruction. Sleep apnea is a more serious condition characterized by repeated episodes of complete or partial airway collapse during sleep, leading to intermittent hypoxia, frequent arousal, and fragmented sleep. Obstructive sleep apnea, the most common form, is associated with obesity, hypertension, and cardiovascular disease; treatment includes continuous positive airway pressure therapy, oral appliances, weight loss, and positional therapy.

\medterm{Solitary Pulmonary Nodule}-- Single, well-circumscribed lesion < 3 cm surrounded by normal lung. Differential includes granuloma, hamartoma, and carcinoma. PET-CT and tissue sampling help determine etiology.

\medterm{Spontaneous Breathing Trial}-- A period of unsupported breathing to assess a patient's readiness for extubation. The patient breathes through a T-piece or with minimal pressure support for 30--120 minutes. Successful SBT predicts successful extubation.

\medterm{Spontaneous Pneumothorax}-- Pneumothorax without trauma. Primary in healthy individuals; secondary with underlying lung disease. Treatment depends on size.

\medterm{Staphylococcal Pneumonia}-- Severe pneumonia from S. aureus, often post-influenza. May cause necrotizing pneumonia, lung abscess, and empyema.

\medterm{Stridor}-- A high-pitched, monophonic, adventitious sound heard primarily during inspiration, caused by narrowing of the extrathoracic or proximal intrathoracic airway. It is a medical emergency when severe. Causes include foreign body aspiration, croup, epiglottitis, and tracheal stenosis.

\medterm{SVC Syndrome}-- Obstruction of the superior vena cava causing facial swelling, plethora, and distended neck veins. Commonly from lung cancer (SCLC, squamous) and lymphoma. Oncologic emergency.

\medterm{Tactile Fremitus}-- Vibrations transmitted through the lung to the chest wall, felt during palpation. It is increased in consolidation (pneumonia) due to better sound transmission through solid tissue. It is decreased in pleural effusion, pneumothorax, and atelectasis due to impaired transmission.

\medterm{Tension Pneumothorax}-- See emergency\_medicine.

\medterm{Thoracentesis}-- Removal of pleural fluid via needle insertion through the chest wall. Performed for diagnosis and symptom relief. Ultrasound guidance improves safety.

\medterm{Thrombolysis for PE}-- Systemic thrombolytics (alteplase) for massive PE with hemodynamic instability. Rapid clot dissolution but significant bleeding risk.

\medterm{Total Lung Capacity}-- Total lung air volume after maximal inspiration. Measured by body plethysmography. Increased in obstructive disease, decreased in restrictive disease. TLC < 80\% confirms restriction.

\medterm{Tracheal Foreign Body}-- A foreign object lodged in the trachea, causing severe respiratory distress, stridor, and complete airway obstruction. It is a medical emergency requiring immediate bronchoscopic removal. In complete obstruction, emergency cricothyrotomy may be needed.

\medterm{Tracheitis}-- Inflammation of the tracheal mucosa. Acute bacterial tracheitis: a rare but life-threatening cause of upper airway obstruction in children, typically presenting as a superinfection following viral croup (parainfluenza). Caused by Staphylococcus aureus (most common), Streptococcus pneumoniae, Haemophilus influenzae type b (now rare due to Hib vaccine), and Moraxella catarrhalis. Presents with high fever, stridor (biphasic, worsening), barky cough, hoarseness, respiratory distress, toxic appearance, and copious purulent tracheal secretions. Diagnosis: lateral neck X-ray—tracheal narrowing ('steeple sign' similar to croup, but with irregular tracheal margins and pseudomembrane), bronchoscopy (definitive—shows inflammation, ulceration, and thick purulent secretions/ pseudomembranes). Treatment: hospitalisation, IV antibiotics (anti-staphylococcal—flucloxacillin or cefuroxime, plus metronidazole for anaerobes), humidified oxygen, and intubation if severe respiratory distress (airway protection, suction of thick secretions). Endotracheal intubation may be required for 5--7 days. Complications: airway obstruction, pneumonia, septic shock, and post-intubation stenosis.

\medterm{Tracheobronchomalacia}-- Abnormal softening of the trachea and/or bronchi causing excessive airway collapse during breathing. It may be congenital or acquired. Symptoms include chronic cough, recurrent infections, and expiratory airflow obstruction. Diagnosis is by dynamic CT or bronchoscopy. Treatment includes CPAP, stenting, or surgical reconstruction.

\medterm{Tracheomalacia}-- Softening of the tracheal cartilage causing excessive collapse of the trachea during expiration. It may be congenital or acquired (prolonged intubation, relapsing polychondritis). Symptoms include expiratory stridor and difficulty clearing secretions. Diagnosis is by dynamic CT or bronchoscopy. Treatment may require tracheal stenting or surgery.

\medterm{Tracheostomy}-- Surgical creation of an airway opening in the anterior neck into the trachea. Indications include prolonged ventilation (> 10--14 days), upper airway obstruction, and airway protection.

\medterm{Transudative Pleural Effusion}-- From systemic factors altering hydrostatic/oncotic pressures. Most commonly from heart failure. Low protein, low LDH. Treatment addresses the underlying cause.

\medterm{Trepopnea}-- Dyspnea that occurs in one lateral decubitus position but not the other. It is typically worse when lying on the side of the affected lung. It may be caused by unilateral lung disease, pleural effusion, or mediastinal mass.

\medterm{Upper Airway Obstruction}-- Narrowing of the airway above the carina, causing inspiratory stridor, dyspnea, and characteristic changes on flow-volume loops (flattening of the inspiratory limb). Causes include tumors, stenosis, goiter, and vocal cord dysfunction. It is a medical emergency when severe.

\medterm{Ventilator-Associated Pneumonia}-- Pneumonia > 48 hours after intubation. Most common ICU-acquired infection. Diagnosis uses quantitative BAL or protected specimen brush cultures.

\medterm{Vocal Resonance}-- The transmission of vocal sounds to the chest wall, assessed by having the patient speak while auscultating. It is increased in consolidation and decreased in pleural effusion or pneumothorax. It correlates with bronchophony and whispered pectoriloquy.

\medterm{Volutrauma}-- Lung injury from excessive tidal volume causing alveolar overdistension. Prevented by low tidal volume ventilation (6 mL/kg IBW). Key component of ARDS management.

\medterm{Westermark Sign}-- Focal area of oligemia (decreased vascularity) distal to an occluded pulmonary artery segment on chest radiography in acute Pulmonary Embolism. \textit{See Pulmonary Embolism (PE).}

\medterm{Whispered Pectoriloquy}-- The ability to hear whispered words clearly through the stethoscope over consolidated lung tissue. It is a sign of pulmonary consolidation. Normally, whispered words are barely audible over the lung fields. It is one of the classic signs of pneumonia on physical examination.

\medterm{Windpipe}-- The windpipe, a 10-12 cm fibrocartilaginous tube from the larynx (C6) to the carina (T4/T5) where it bifurcates into main bronchi. See \hyperlink{Trachea}{Trachea}.

\medterm{Pulmonary Hypertension}-- Abnormally increased pressure in the pulmonary circulation, confirmed by right-heart catheterization when mean pulmonary artery pressure exceeds 20 mmHg at rest. Classification by cause guides treatment and includes pulmonary arterial, left-heart, lung-disease/hypoxia, chronic thromboembolic, and multifactorial groups.

\medterm{Pulmonary Arterial Hypertension (PAH)}-- A precapillary pulmonary hypertension syndrome caused by increased pulmonary vascular resistance from pulmonary arteriolar disease. It may be idiopathic, heritable, drug-associated, or associated with conditions such as connective-tissue disease, congenital heart disease, portal hypertension, or HIV.

\medterm{Chronic Thromboembolic Pulmonary Hypertension (CTEPH)}-- Pulmonary hypertension caused by organized, persistent thromboembolic obstruction and secondary pulmonary-vessel remodeling after pulmonary embolism. It is potentially curable in selected patients by pulmonary endarterectomy or balloon pulmonary angioplasty.

\medterm{Acute Hypoxemic Respiratory Failure}-- Inability of the respiratory system to maintain adequate arterial oxygenation, usually manifested by low arterial oxygen or oxygen saturation despite appropriate inspired oxygen. Causes include pneumonia, pulmonary edema, ARDS, atelectasis, and pulmonary embolism.

\medterm{Acute Hypercapnic Respiratory Failure}-- Inability to eliminate carbon dioxide, producing elevated arterial PaCO$_2$ and often respiratory acidemia. Causes include COPD exacerbation, severe asthma, neuromuscular weakness, central respiratory depression, and chest-wall restriction.

\medterm{Respiratory Drive}-- The neural output that stimulates respiratory muscles and determines ventilation. It is generated by brainstem networks and modified by chemoreceptor input, mechanoreceptors, wakefulness, metabolic state, medications, and behavioral control.

\medterm{Respiratory Quotient (RQ)}-- The ratio of carbon dioxide production to oxygen consumption for a particular substrate or metabolic state. It differs from the respiratory exchange ratio measured at the mouth and is used in metabolic and ventilatory physiology.

\medterm{Alveolar Ventilation Equation}-- The relationship between alveolar ventilation and arterial carbon-dioxide tension, approximately PaCO$_2$ = 0.863 times carbon-dioxide production divided by alveolar ventilation. It explains why reduced effective ventilation causes hypercapnia.

\medterm{Bronchodilator Response}-- The change in spirometric airflow after administration of a short-acting bronchodilator. A significant response can support variable airflow limitation, but a negative response does not exclude asthma and a positive response is not specific to one disease.

\medterm{Cystic Fibrosis}-- An autosomal recessive multisystem disorder caused by pathogenic CFTR variants that impair epithelial chloride and bicarbonate transport. Thick secretions cause chronic sinopulmonary infection, bronchiectasis, pancreatic insufficiency, intestinal obstruction, infertility, and abnormal sweat chloride.

\medterm{Sarcoidosis}-- A multisystem inflammatory disease characterized by noncaseating granulomas in affected tissues. It commonly involves the lungs and intrathoracic lymph nodes but may affect skin, eyes, heart, nervous system, and other organs; diagnosis requires compatible findings and exclusion of other granulomatous causes.


\medterm{Acute Exacerbation of Asthma}-- A progressive increase in dyspnea, cough, wheezing, or chest tightness with reduced expiratory airflow compared with the patient’s baseline. Severe attacks may cause exhaustion, silent chest, hypercapnia, respiratory failure, and need for urgent treatment.

\medterm{COPD Exacerbation}-- Acute worsening of respiratory symptoms beyond normal day-to-day variation in a patient with COPD that leads to additional therapy. Common triggers include viral or bacterial infection, air pollution, and cardiopulmonary comorbidity.

\medterm{Pulmonary Rehabilitation}-- A supervised, multidisciplinary program combining exercise training, education, behavioral support, and self-management for chronic respiratory disease. It improves exercise capacity, dyspnea, functional status, and quality of life.

\medterm{Sleep-Related Hypoventilation}-- Inadequate ventilation during sleep causing elevated carbon dioxide, oxygen desaturation, or both. Causes include obesity, neuromuscular disease, chest-wall disorders, chronic lung disease, medications, and central respiratory-control abnormalities.


\medterm{Ventilator-Induced Lung Injury (VILI)}-- Lung injury caused or amplified by mechanical ventilation, including volutrauma, barotrauma, atelectrauma, and biotrauma. Lung-protective ventilation limits tidal volume and excessive airway pressures while maintaining appropriate oxygenation and carbon-dioxide targets.
\medterm{Air bronchogram}-- A radiographic sign showing air-filled bronchi visible against opacified lung parenchyma, indicating that the airways remain patent while the surrounding alveoli are filled with fluid, pus, or cells.
\medterm{Chronic eosinophilic pneumonia}-- An idiopathic inflammatory lung disease characterized by eosinophilic infiltration of the pulmonary parenchyma, presenting with fever, cough, dyspnea, and peripheral airspace opacities on imaging.
\medterm{Eosinophilic bronchitis}-- Inflammation of the bronchial airways with eosinophilic infiltration causing chronic cough, with normal spirometry and no airway hyperresponsiveness, distinguishing it from asthma.
\medterm{Hydropneumothorax}-- The simultaneous presence of both fluid and air in the pleural space, causing respiratory distress and dullness to percussion with absent breath sounds, typically requiring chest tube drainage.
\medterm{Mycoplasma pneumoniae-associated mucositis}-- Inflammatory lesions of the oral, ocular, or genital mucosa occurring as an extrapulmonary manifestation of Mycoplasma pneumoniae infection, often presenting with target-like skin lesions.
\medterm{Oleothorax}-- Historical surgical obliteration of the pleural space using infused oil to collapse the lung, once used for tuberculosis treatment before effective chemotherapy, now largely obsolete.
\medterm{Plastic bronchitis}-- A rare condition producing large, branching mucous casts that form in the airways, causing life-threatening airway obstruction, often associated with congenital heart disease or post-fontan circulation.
\medterm{Pleural perforation}-- A breach in the pleural lining allowing air or fluid to enter the pleural space, potentially causing pneumothorax or hydropneumothorax, often from trauma, procedures, or underlying lung disease.
\medterm{Pneumocephalus}-- The presence of air or gas within the cranial cavity, most commonly following trauma, surgery, or sinus disease, and potentially causing increased intracranial pressure if tension develops.
\medterm{Reexpansion pulmonary edema}-- Noncardiogenic pulmonary edema developing rapidly after rapid reexpansion of a collapsed lung, typically following chest tube drainage of a large or long-standing pneumothorax.
\medterm{Asthma exacerbation}-- An acute or subacute worsening of asthma symptoms and airflow limitation, typically causing increased wheeze, cough, chest tightness, or breathlessness.

\medterm{Hyperventilation Syndrome}-- Recurrent or persistent overbreathing that exceeds metabolic requirements, usually producing reduced arterial carbon dioxide and respiratory alkalosis in the absence of primary pulmonary disease. Symptoms may include breathlessness, chest tightness, dizziness, palpitations, paraesthesia, and carpopedal spasm. It is a diagnosis of exclusion after serious cardiopulmonary and metabolic causes have been assessed.

\medterm{Bronchiolitis Obliterans}-- An uncommon, often irreversible obstructive small-airway disease caused by inflammation and fibrotic narrowing or obliteration of the bronchioles. Causes include severe infection, inhalational injury, autoimmune disease, and complications after lung or bone-marrow transplantation. It presents with progressive breathlessness, dry cough, wheezing, and fixed airflow obstruction; diagnosis relies on clinical history, pulmonary-function testing, and chest CT.

\medterm{Acute Bronchitis}-- Acute inflammation of the tracheobronchial tree, usually caused by a viral respiratory infection and characterized by cough, with or without sputum production, lasting up to several weeks. Wheeze, chest discomfort, and low-grade fever may occur; persistent high fever, hypoxaemia, focal findings, or severe illness should prompt assessment for pneumonia or another cause.

\medterm{Apnea}-- See \hyperlink{Apnoea}{Apnoea}.

\medterm{Clubbing}-- Bulbous enlargement of the distal fingers or toes with increased nail curvature and loss of the normal nail-fold angle. It reflects proliferation of connective tissue in the distal digits and may be associated with chronic pulmonary disease, cyanotic heart disease, inflammatory bowel disease, or malignancy.

\medterm{Hypercapnia}-- An abnormally elevated arterial partial pressure of carbon dioxide, usually caused by alveolar hypoventilation. It may produce headache, somnolence, confusion, asterixis, respiratory acidosis, and, when severe or rapidly developing, coma.

\medterm{Intercostal Retraction}-- Inward movement of the skin and soft tissues between the ribs during inspiration, indicating increased work of breathing. It is common in infants and children with respiratory distress and may accompany airway obstruction, bronchiolitis, pneumonia, asthma, or other pulmonary disease.

\medterm{Lymphangiectasis}-- Abnormal dilatation of lymphatic vessels. It may be congenital or acquired and can cause lymph leakage, oedema, protein-losing enteropathy, chylous effusions, or pulmonary and pleural complications depending on the affected site.

\medterm{Radiation Pneumonitis}-- An inflammatory lung injury occurring weeks to months after thoracic radiation. It may cause dry cough, fever, dyspnea, and hypoxemia within the irradiated field and can progress to radiation fibrosis.

\medterm{Tracheobronchomegaly}-- Abnormal enlargement of the trachea and main bronchi, usually associated with connective-tissue weakness. It may cause recurrent lower-respiratory infection, bronchiectasis, ineffective cough, dyspnea, and difficulty clearing airway secretions.

\medterm{Chemical Pneumonitis}-- Inflammation and lung injury caused by inhalation or aspiration of irritating chemicals, gastric contents, hydrocarbons, or other toxic substances. It may cause cough, dyspnea, hypoxemia, wheezing, and diffuse pulmonary infiltrates.

\medterm{Exercise-Induced Asthma}-- Bronchoconstriction occurring during or shortly after physical exertion, causing cough, wheezing, chest tightness, or shortness of breath. It is associated with airway hyperresponsiveness and may occur in people with or without persistent asthma.
